%====================================================================
%  TidalLanes_Model_v1.tex
%
%  Self-contained Model and Estimation write-up for the
%  "Directional Congestion and Tidal Lanes" project.
%
%  Structure
%    Section 1  Model
%    Section 2  Estimation
%    Appendix   Derivations, hat-algebra, and IV strategy
%
%  This file is meant to be compiled stand-alone AND to be \input
%  into the main manuscript by stripping the preamble and the
%  \begin{document}/\end{document} pair.
%====================================================================
\documentclass[12pt]{article}

\usepackage[margin=1in]{geometry}
\usepackage{amsmath,amssymb,amsthm,mathtools}
\usepackage{graphicx,booktabs,tabularx,array,setspace,float,enumitem}
\usepackage[colorlinks=true,linkcolor=blue,citecolor=blue,urlcolor=blue]{hyperref}
\usepackage{comment}
\setstretch{1.15}
\setlength{\parindent}{2em}

\newtheorem{proposition}{Proposition}
\newtheorem{lemma}{Lemma}
\newtheorem{definition}{Definition}
\newtheorem{assumption}{Assumption}
\newtheorem{remark}{Remark}

\DeclareMathOperator*{\argmax}{arg\,max}

% Shorthand
\newcommand{\N}{\mathcal N}
\newcommand{\E}{\mathcal E}
\newcommand{\R}{\mathcal R}
\newcommand{\Lbar}{\bar L}
\newcommand{\Wbar}{\bar W}
\newcommand{\ubar}{\bar u}
\newcommand{\Abar}{\bar A}
\newcommand{\tbar}{\bar t}
\newcommand{\tkl}{t_{kl}}
\newcommand{\tijk}{\tau_{ij}}
\newcommand{\free}{\mathrm{ff}}
\newcommand{\obs}{\mathrm{obs}}
\newcommand{\cf}{\mathrm{cf}}
\newcommand{\spr}{\rho}

\title{Directional Congestion and Tidal Lanes:\\
Model and Estimation Strategy (v1)}
\author{Yige Duan \and Xiaoruo Feng \and Yizhen Gu}
\date{\today}

\begin{document}
\maketitle

\begin{abstract}
\noindent
This companion note develops a complete statement of the quantitative
spatial-equilibrium model and the corresponding estimation strategy used
in ``Directional Congestion and Tidal Lanes''.  Relative to
\cite{AllenArkolakis2022} the model retains the Fr\'echet route-choice
structure and the link-traffic gravity representation but allows all
link primitives (lane capacity, free-flow speed, realized travel time)
and bilateral commuting costs to be \emph{directional}.  The estimation
pipeline is adapted to our data, which contain high-frequency directed
speeds and a grid-level origin--destination matrix but no directly
observed link flows.  Every step is fully derived; practical
implementation choices (calibration vs.\ internal estimation, the IV
strategy for the congestion elasticity, and the exact-hat-algebra
counterfactual solver) are stated as explicit algorithms.
\end{abstract}

\tableofcontents
\clearpage

%====================================================================
\section{Model}\label{sec:model}
%====================================================================

\subsection{Environment and Notation}\label{ssec:env}

\paragraph{Locations and directed network.}
The city is partitioned into a finite set of spatial cells
(\emph{grid cells}) $\N=\{1,\dots,N\}$.  Each node $i\in\N$ has a
geographic centroid and is associated with an area of land.  The road
network is represented as a directed graph $(\N,\E)$ with edge set
$\E\subseteq\N\times\N$.  A directed edge $(k,l)\in\E$ describes an
admissible move from node $k$ to node $l$ along the observed road
infrastructure; $(k,l)\in\E$ does not imply $(l,k)\in\E$, and even when
both hold the two directions are treated as distinct objects with
potentially different primitives.

\paragraph{Directed link primitives.}
Each directed link $(k,l)\in\E$ is characterized by
\begin{itemize}[leftmargin=2em,itemsep=0.2em]
  \item a physical length $\ell_{kl}>0$ (centroid-to-centroid distance),
  \item a directed lane capacity $n_{kl}>0$,
  \item a free-flow speed $v^{0}_{kl}>0$,
\end{itemize}
from which the directed free-flow travel time is
\begin{equation}
  \tbar_{kl}\;\equiv\;\frac{\ell_{kl}}{v^{0}_{kl}}.
  \label{eq:tbar_def}
\end{equation}
The realized, congestion-inclusive directed travel time is denoted
$\tkl>0$ and is determined in equilibrium by \S\ref{ssec:cong}.
Directional asymmetry is a first-order feature of the environment:
\[
  \ell_{kl}\neq \ell_{lk},\qquad
  n_{kl}\neq n_{lk},\qquad
  v^{0}_{kl}\neq v^{0}_{lk},\qquad
  \tkl\neq t_{lk},
\]
are all permitted in general.

\paragraph{Population and fundamentals.}
The city is populated by a unit mass of workers of total measure
$\Lbar$, allocated across residences $\{L_i^R\}_{i\in\N}$ and
workplaces $\{L_j^F\}_{j\in\N}$, with
\begin{equation}
  \sum_{i\in\N} L_i^R
  \;=\;
  \sum_{j\in\N} L_j^F
  \;=\;
  \Lbar.
  \label{eq:pop}
\end{equation}
Each location is endowed with a baseline residential amenity
$\ubar_i>0$ and a baseline workplace productivity $\Abar_j>0$.
Endogenous amenity and productivity depend on the local scale of
activity through
\begin{equation}
  u_i=\ubar_i\bigl(L_i^R\bigr)^{\beta},
  \qquad
  A_j=\Abar_j\bigl(L_j^F\bigr)^{\alpha},
  \label{eq:fundamentals}
\end{equation}
where $\alpha$ is the \emph{production-externality (agglomeration)
elasticity} and $\beta$ is the \emph{residential-externality
(congestion) elasticity}, following the convention of
\cite{AllenArkolakis2014,AllenArkolakis2022}.  We impose the sign
restrictions
\begin{equation}
  \alpha\;>\;0,
  \qquad
  \beta\;<\;0,
  \label{eq:sign}
\end{equation}
so that employment density in a workplace cell raises productivity
through agglomeration spillovers (knowledge spillovers, thick labour
markets, input--output sharing), while residential density in a
residence cell lowers per-capita amenity through congestion forces
(housing competition, crowding externalities).  

\begin{comment}
	This one-parameter
spillover formulation treats $\alpha$ and $\beta$ as \emph{composite}
elasticities that net out the underlying agglomeration and congestion
forces; it is the standard specification in the Allen--Arkolakis
class and is distinct from the explicit two-force decomposition used
by \cite{Ahlfeldt2015} (see Remark~\ref{rem:AA_vs_Ahlfeldt}).
\end{comment}


\paragraph{Parameters.}
The model is parameterized by
\begin{itemize}[leftmargin=2em,itemsep=0.2em]
  \item $\theta>0$, the Fr\'echet shape parameter governing
        heterogeneity in route (and residence--workplace) choice;
  \item $\alpha,\beta\in\mathbb R$, the externality elasticities in
        \eqref{eq:fundamentals};
  \item $\lambda\ge0$, the congestion elasticity that maps traffic per
        lane into realized travel time (\S\ref{ssec:cong}).
\end{itemize}
In our baseline implementation $(\theta,\alpha,\beta)$ are externally
calibrated and $\lambda$ is estimated from a reduced-form speed
regression (Section~\ref{sec:est}).

\begin{remark}[AA vs.\ Ahlfeldt parameterization of spillovers]
\label{rem:AA_vs_Ahlfeldt}
\cite{Ahlfeldt2015} model production and residential externalities as
the explicit sum of an \emph{agglomeration force} (with elasticity
$\lambda^A>0$ and spatial decay $\delta^A>0$) and an offsetting
\emph{congestion force} (with elasticity $\lambda^C<0$ and decay
$\delta^C$).  Their composite, reduced-form elasticities are obtained
from these primitive forces and can in principle be of either sign
depending on which force dominates locally.  By contrast, the
Allen--Arkolakis class adopts a one-parameter, sign-restricted
specification in which $\alpha>0$ summarises net agglomeration on the
production side and $\beta<0$ summarises net congestion on the
residential side.  Our setting follows AA: identification of the
underlying primitive forces $(\lambda^A,\lambda^C,\delta^A,\delta^C)$
would require additional moment conditions (e.g., spatially-decaying
density measures) that our Baidu commuting matrix does not deliver,
so we work directly with the composite $(\alpha,\beta)$.
\end{remark}

\begin{remark}[Why directionality matters]
In standard monocentric commuting models the bilateral cost matrix
$\{\tijk\}$ is symmetric and a lane reallocation that preserves the
total capacity $n_{kl}+n_{lk}$ is a pure re-labelling.  Once travel
times are directional, however, a reallocation of lanes from the
off-peak direction to the peak direction strictly lowers $\tkl$ and
strictly raises $t_{lk}$; the induced change in $\{\tijk\}$ is
asymmetric across origin--destination pairs and has real effects on
flows, location choice, and welfare.  This is the channel through
which tidal-lane policies operate in our model.
\end{remark}

\subsection{Individuals, Routes, and Transport Cost}\label{ssec:ind}

\paragraph{Choice problem.}
A representative worker chooses simultaneously a residence $i\in\N$,
a workplace $j\in\N$, and a commuting route $r\in\R_{ij}$.  A route is
a finite ordered sequence of nodes
\[
  r=(r_{0},r_{1},\dots,r_{K})
  \quad\text{with}\quad
  r_{0}=i,\;\; r_{K}=j,\;\; (r_{m-1},r_{m})\in\E\ \forall m=1,\dots,K,
\]
where $K=K(r)$ is the (finite) length of the route.  The set of
feasible routes $\R_{ij}$ is countably infinite when cycles are
admitted.  The indirect utility of choice $(i,j,r)$ is
\begin{equation}
  V_{ij,r}
  \;=\;
  \left(
    \frac{u_i\,w_j}
         {\prod_{m=1}^{K(r)} t_{r_{m-1},r_{m}}}
  \right)\varepsilon_{ij,r},
  \qquad
  \varepsilon_{ij,r}\stackrel{\mathrm{iid}}{\sim}\mathrm{Fr\acute echet}(\theta),
  \label{eq:utility}
\end{equation}
where $w_j$ is the wage at workplace $j$ and $\varepsilon_{ij,r}$ is
an idiosyncratic preference shock drawn independently across
$(i,j,r)$-triples from a Fr\'echet distribution with shape $\theta$
and unit scale, i.e.\ $\Pr(\varepsilon\le x)=\exp(-x^{-\theta})$
for $x>0$.

\begin{remark}[Time units and the $t_{kl}$ convention]\label{rem:iceberg}
We treat $\tkl=\ell_{kl}/v_{kl}$ as directed physical travel time,
entering utility \eqref{eq:utility} and link-traffic gravity through
the same multiplicative (``iceberg-style'') channel: the product
$\prod_{m=1}^{K(r)}t_{r_{m-1},r_{m}}$ is the cumulative travel-cost
factor on route~$r$.  The congestion technology
\eqref{eq:cong_primitive} is therefore a primitive on physical travel
time, and the speed-flow regression of Section~\ref{ssec:lambdahat}
recovers $\lambda$ without a $\theta$-rescaling: the elasticity of
$\tkl$ with respect to traffic per lane equals minus the elasticity
of $v_{kl}$ with respect to traffic per lane.

\cite{AllenArkolakis2022} work with the transformed iceberg index
$\tilde t_{kl}=\tkl^{1/\theta}$, equivalently with $\delta_0=1/\theta$
in the micro-foundation
$\tilde t_{kl}=\bigl(\ell_{kl}/v_{kl}\bigr)^{\delta_0}$ and
$v_{kl}=v^{0}_{kl}(n_{kl}/\xi_{kl})^{\lambda/\delta_0}$.  That
transformation renormalises the Fr\'echet bilateral cost index
$\tijk^{-\theta}$ so that its distance elasticity equals $-1$ rather
than $-\theta$.  The two conventions are algebraically equivalent up
to a rescaling of the congestion elasticity,
$\lambda^{\mathrm{AA}}=\lambda/\theta$, and deliver identical
equilibrium predictions.  The physical-time convention is adopted
here because it keeps the mapping from the speed-flow regression
coefficient to $\lambda$ direct, without a factor-of-$\theta$
rescaling; conversion to the AA scaling is a one-line adjustment
wherever the alternative is useful for cross-paper comparison.

A change of the base time unit (seconds to minutes, minutes to
hours) rescales every $\tkl$ by a common factor that is absorbed
into $\chi$ in \eqref{eq:chi_def} and leaves the equilibrium
distribution $\{L_i^R,L_j^F,\pi_{ij},\xi_{kl}\}$ unchanged.
\end{remark}

\paragraph{Choice probabilities.}
Standard properties of i.i.d.\ Fr\'echet draws imply that the
probability of choosing $(i,j,r)$ is
\begin{equation}
  \pi_{ij,r}
  \;=\;
  \frac{
    \bigl(u_i\,w_j\,\prod_{m=1}^{K(r)} t_{r_{m-1},r_{m}}^{-1}\bigr)^{\theta}
  }{
    \displaystyle\sum_{(i',j')\in\N^{2}}
    \sum_{r'\in\R_{i'j'}}
    \bigl(u_{i'}\,w_{j'}\,\prod_{m=1}^{K(r')} t_{r'_{m-1},r'_{m}}^{-1}\bigr)^{\theta}
  }.
  \label{eq:routeprob}
\end{equation}

\paragraph{Transport cost index.}
The quantity $u_i\,w_j$ in the numerator of \eqref{eq:routeprob} does
not depend on the route.  Summing over routes $r\in\R_{ij}$
motivates the definition of the \emph{bilateral commuting cost
index} $\tijk$ via
\begin{equation}
  \tijk^{-\theta}
  \;\equiv\;
  \sum_{r\in\R_{ij}}
  \prod_{m=1}^{K(r)} t_{r_{m-1},r_{m}}^{-\theta}.
  \label{eq:tau_def}
\end{equation}
This object summarizes \emph{all} feasible routes from $i$ to $j$ and
inherits any directional asymmetry present in the underlying link
times: in general $\tijk\neq\tau_{ji}$.

\paragraph{Series representation.}
Define the directed weighted adjacency matrix $A=[a_{kl}]$ with
\begin{equation}
  a_{kl}\;\equiv\;\tkl^{-\theta}\;\text{if}\;(k,l)\in\E
  \quad\text{and}\quad a_{kl}\equiv 0\;\text{otherwise}.
  \label{eq:Adef}
\end{equation}
For any integer $K\ge 0$, $(A^{K})_{ij}$ equals the sum of
$\prod_{m=1}^{K} a_{r_{m-1},r_{m}}$ taken over all directed walks of
exactly $K$ edges from $i$ to $j$.  Summing over walk lengths and
using \eqref{eq:tau_def} delivers the series representation
$\tijk^{-\theta}=\sum_{K\ge 0}(A^{K})_{ij}$; see
Appendix~\ref{app:tau} for details.

\begin{assumption}[Spectral radius condition]\label{ass:specrad}
The matrix $A$ defined in \eqref{eq:Adef} satisfies $\spr(A)<1$,
where $\spr(\cdot)$ denotes the spectral radius.
\end{assumption}

\begin{proposition}[Network representation]\label{prop:tau_series}
Under Assumption~\ref{ass:specrad}, $\tijk$ in \eqref{eq:tau_def} is
finite for every $(i,j)\in\N^2$ and
\begin{equation}
  \tijk^{-\theta}
  \;=\;
  \sum_{K=0}^{\infty}(A^{K})_{ij}
  \;=\;\bigl[(I-A)^{-1}\bigr]_{ij}
  \;\equiv\; b_{ij},
  \qquad
  \tijk=b_{ij}^{-1/\theta}.
  \label{eq:tau_series}
\end{equation}
\end{proposition}

The proof is given in Appendix~\ref{app:tau} and follows the
combinatorial interpretation of matrix powers combined with the
geometric-series bound.  Because $A$ is not symmetric, neither is
$(I-A)^{-1}$, and therefore $\tijk\neq\tau_{ji}$ whenever the
underlying link times are directional.

\begin{remark}[Finiteness of the cost index under arbitrary cycles]
\label{rem:cycles}
The series in \eqref{eq:tau_series} includes all walks, including
those that revisit nodes and edges.  Assumption~\ref{ass:specrad} is
the sharp condition for the series to converge.  Because
$a_{kl}=\tkl^{-\theta}$ is decreasing in $\tkl$ and $\theta$ is large,
the assumption is satisfied in all of our numerical implementations
on grid-level graphs at AM-peak speeds.
\end{remark}

\subsection{Firms}\label{ssec:firms}

Firms at each workplace $j\in\N$ produce a single homogeneous
numeraire good using labor as the sole input and operate under
perfect competition.  Labor productivity equals $A_j$ in
\eqref{eq:fundamentals} and the market-clearing wage is
\begin{equation}
  w_j\;=\;A_j\;=\;\Abar_j\bigl(L_j^F\bigr)^{\alpha}.
  \label{eq:wage}
\end{equation}
Because we focus on commuting and transportation, we do not model
goods trade, housing prices, or non-traded amenities separately;
$u_i$ in \eqref{eq:fundamentals} is interpreted as a composite
residential amenity that already absorbs the equilibrium effect of
housing market clearing and is consistent with the reduced-form
parameterization in \cite{AllenArkolakis2022}.

\subsection{Commuting Flows: Gravity Equation}\label{ssec:gravity}

We now aggregate the route-choice probabilities \eqref{eq:routeprob}
to obtain the gravity equation for commuting flows.  The derivation
proceeds in three steps: (i) re-writing each $V_{ij,r}$ as a
multiplicative ``scale $\times$ shock'' expression, (ii) using the
max-stability of the Fr\'echet distribution to compute the
distribution of the worker's overall maximum utility, and
(iii) identifying expected welfare with a single scalar index
$\Wbar$.

\paragraph{Scale factors.}
For each triple $(i,j,r)$ define the \emph{scale factor}
\begin{equation}
  \omega_{ij,r}
  \;\equiv\;
  u_i\,w_j\,\prod_{m=1}^{K(r)} t_{r_{m-1},r_{m}}^{-1},
  \qquad
  V_{ij,r}=\omega_{ij,r}\,\varepsilon_{ij,r}.
  \label{eq:omega_def}
\end{equation}
Because $\varepsilon_{ij,r}$ is Fr\'echet with shape $\theta$ and
unit scale, $\omega_{ij,r}\varepsilon_{ij,r}$ is Fr\'echet with
shape $\theta$ and scale $\omega_{ij,r}$.

\paragraph{Maximum across $(i,j,r)$.}
Let $M\equiv\max_{(i,j,r)} V_{ij,r}$ be the worker's overall maximum
utility.  By independence across $(i,j,r)$,
\begin{align}
  \Pr(M\le y)
  &=\prod_{(i,j,r)}\Pr(V_{ij,r}\le y)
   =\prod_{(i,j,r)}\exp\!\left(-(y/\omega_{ij,r})^{-\theta}\right)\nonumber\\
  &=\exp\!\left(-y^{-\theta}\sum_{(i,j,r)}\omega_{ij,r}^{\theta}\right)
   =\exp\!\left(-y^{-\theta}S\right),\label{eq:Fmax}
\end{align}
where
\begin{equation}
  S\;\equiv\;\sum_{(i,j)\in\N^{2}}\sum_{r\in\R_{ij}}\omega_{ij,r}^{\theta}
  \;=\;\sum_{(i,j)\in\N^{2}} u_i^{\theta}w_j^{\theta}\,\tijk^{-\theta},
  \label{eq:S_def}
\end{equation}
and the last equality uses the definition of $\tijk$ in
\eqref{eq:tau_def}.  Equation~\eqref{eq:Fmax} establishes that $M$ is
itself Fr\'echet with shape $\theta$ and scale $S^{1/\theta}$, so
\begin{equation}
  \mathbb E[M]=\Gamma\!\left(1-\tfrac{1}{\theta}\right)S^{1/\theta},
  \qquad
  \Wbar\;\equiv\;S^{1/\theta}
  =\left(\sum_{(i,j)}\tijk^{-\theta}u_i^{\theta}w_j^{\theta}\right)^{\!1/\theta}.
  \label{eq:W_def}
\end{equation}
We identify $\Wbar$ with (Fr\'echet-scale-normalized) expected
welfare.  See Appendix~\ref{app:welfare} for the step-by-step
derivation.

\paragraph{Choice probabilities at the $(i,j)$ level and gravity.}
The probability that a worker chooses the residence--workplace pair
$(i,j)$, summed across routes, equals
\begin{equation}
  \Pr(i,j)
  \;=\;
  \sum_{r\in\R_{ij}}\pi_{ij,r}
  \;=\;
  \frac{u_i^{\theta}w_j^{\theta}\tijk^{-\theta}}{S}.
  \label{eq:probij}
\end{equation}
Multiplying by the total population $\Lbar$ yields the \emph{gravity
equation} for commuting flows:
\begin{equation}
  \boxed{\;L_{ij}
  \;=\;
  \Lbar\cdot\Pr(i,j)
  \;=\;
  \frac{\tijk^{-\theta}\,u_i^{\theta}\,w_j^{\theta}}{\Wbar^{\theta}}\Lbar.\;}
  \label{eq:gravity}
\end{equation}
Marginal populations satisfy the market-clearing conditions
\begin{equation}
  L_i^R\;=\;\sum_{j\in\N}L_{ij},
  \qquad
  L_j^F\;=\;\sum_{i\in\N}L_{ij},
  \qquad
  \sum_{i}L_i^R=\sum_{j}L_j^F=\Lbar.
  \label{eq:marketclearing}
\end{equation}

\subsection{Market Access and Link Traffic}\label{ssec:ma}

\paragraph{Market-access representation of gravity.}
Summing the gravity equation \eqref{eq:gravity} over $j$ and using
$L_i^R=\sum_j L_{ij}$ gives
\begin{equation}
  L_i^R
  \;=\;\frac{u_i^{\theta}\Lbar}{\Wbar^{\theta}}
       \sum_{j\in\N}\tijk^{-\theta}w_j^{\theta}.
  \label{eq:Li_expanded}
\end{equation}
Symmetrically, summing over $i$ gives $L_j^F$.  Motivated by
\eqref{eq:Li_expanded} we define the \emph{resident market-access
index} $\Pi_i$ and the \emph{firm market-access index} $P_j$ by
\begin{align}
  \Pi_i^{-\theta}
  &\;\equiv\;
  \sum_{j\in\N}\tijk^{-\theta}\,L_j^F\,P_j^{\theta},
  \label{eq:Pi_def}\\[4pt]
  P_j^{-\theta}
  &\;\equiv\;
  \sum_{i\in\N}\tijk^{-\theta}\,L_i^R\,\Pi_i^{\theta}.
  \label{eq:Pj_def}
\end{align}
These are the asymmetric analogues of the resident and firm
commuter-market-access aggregators in the Ahlfeldt et al.\
\cite{Ahlfeldt2015} tradition, where the bilateral cost term
$\tijk^{-\theta}$ is replaced by its directional counterpart.

\paragraph{Market-access form of gravity.}
Combining \eqref{eq:Pi_def}--\eqref{eq:Pj_def} with the gravity
equation \eqref{eq:gravity} yields the equivalent representation
\begin{equation}
  L_{ij}
  \;=\;
  \tijk^{-\theta}\cdot
  \frac{L_i^R}{\Pi_i^{-\theta}}\cdot
  \frac{L_j^F}{P_j^{-\theta}}.
  \label{eq:gravity_MA}
\end{equation}
Equation~\eqref{eq:gravity_MA} is useful because it expresses
bilateral flows in terms of observable marginal populations
$(L_i^R,L_j^F)$ and two scalar accessibility indices per node, and
because it is the form we will use to derive link-level traffic.

\paragraph{Link-traversal probabilities.}
For any directed edge $(k,l)\in\E$ and any origin--destination pair
$(i,j)$, let $\pi^{kl}_{ij}$ denote the expected number of times a
worker commuting from $i$ to $j$ traverses edge $(k,l)$ along her
chosen route.\footnote{Because our route set $\R_{ij}$ includes walks
with repeated edges, $\pi^{kl}_{ij}$ should be interpreted as an
expected count rather than a probability bounded by one; this is
immaterial for the aggregate traffic computation below.}
Routes that traverse $(k,l)$ can be decomposed as a walk from $i$ to
$k$, followed by the single edge $(k,l)$, followed by a walk from
$l$ to $j$.  Summing the product of the corresponding walk weights
and dividing by $\tijk^{-\theta}$ gives
\begin{equation}
  \pi^{kl}_{ij}
  \;=\;\frac{\tau_{ik}^{-\theta}\,\tkl^{-\theta}\,\tau_{lj}^{-\theta}}
            {\tijk^{-\theta}}
  \;=\;\left(\frac{\tijk}{\tau_{ik}\,\tkl\,\tau_{lj}}\right)^{\!\theta}.
  \label{eq:piijkl}
\end{equation}
The detailed derivation is in Appendix~\ref{app:linkgravity}.

\paragraph{Directed link traffic $\xi_{kl}$.}
Aggregating \eqref{eq:piijkl} across origin--destination pairs with
the gravity flows \eqref{eq:gravity_MA} gives the directed traffic
on edge $(k,l)$:
\begin{equation}
  \xi_{kl}
  \;\equiv\;\sum_{i\in\N}\sum_{j\in\N}L_{ij}\,\pi^{kl}_{ij}.
  \label{eq:xi_def}
\end{equation}
The following key result has a compact multiplicative form.

\begin{proposition}[Link-traffic gravity]\label{prop:linkgravity}
Under Assumption~\ref{ass:specrad} and definitions
\eqref{eq:Pi_def}--\eqref{eq:Pj_def},
\begin{equation}
  \boxed{\;\xi_{kl}\;=\;\tkl^{-\theta}\,P_k^{-\theta}\,\Pi_l^{-\theta}.\;}
  \label{eq:xi_gravity}
\end{equation}
\end{proposition}
\begin{proof}[Sketch]
Substitute \eqref{eq:gravity_MA} and \eqref{eq:piijkl} into
\eqref{eq:xi_def}.  The numerator factorizes as
$\tau_{ik}^{-\theta}\tkl^{-\theta}\tau_{lj}^{-\theta}\cdot
L_i^R\Pi_i^{\theta}\cdot L_j^F P_j^{\theta}$.  Separating sums over
$i$ and over $j$ and invoking \eqref{eq:Pi_def}--\eqref{eq:Pj_def}
yields $\tkl^{-\theta}P_k^{-\theta}\Pi_l^{-\theta}$.  The full
argument appears in Appendix~\ref{app:linkgravity}.
\end{proof}

\begin{remark}[Interpretation]
Equation~\eqref{eq:xi_gravity} is the directed analogue of the
Allen--Arkolakis link-traffic gravity formula.  The three multiplicative
factors admit clean interpretations: $\tkl^{-\theta}$ is the
edge-specific resistance cost, $P_k^{-\theta}$ aggregates all origin
sides that feed into $k$, and $\Pi_l^{-\theta}$ aggregates all
destination sides that can be reached from $l$.  Because every
component is directional, $\xi_{kl}\ne\xi_{lk}$ in general.
\end{remark}

\subsection{Congestion Technology}\label{ssec:cong}

\paragraph{Primitive equation.}
Following the structure of \cite{AllenArkolakis2022}, we assume that
the realized directed travel time on link $(k,l)$ is increasing in
directed traffic per directed lane:
\begin{equation}
  \tkl
  \;=\;
  \tbar_{kl}\!\left(\frac{\xi_{kl}}{n_{kl}}\right)^{\!\lambda},
  \qquad \lambda\ge 0,
  \label{eq:cong_primitive}
\end{equation}
where $\tbar_{kl}=\ell_{kl}/v^{0}_{kl}$ is the directed free-flow
travel time defined in \eqref{eq:tbar_def} and $n_{kl}$ is the
directed lane capacity of the link.  When $\lambda=0$, $\tkl=\tbar_{kl}$
for every $(k,l)$ and the model collapses to a frictionless
network.

\paragraph{Reduced-form representation.}
Substituting the link-traffic gravity equation \eqref{eq:xi_gravity}
into \eqref{eq:cong_primitive} and solving for $\tkl$ yields a
reduced-form expression for equilibrium travel time that holds node
accessibility fixed:
\begin{equation}
  \tkl
  \;=\;
  \tbar_{kl}^{1/(1+\theta\lambda)}\cdot
  n_{kl}^{-\lambda/(1+\theta\lambda)}\cdot
  \bigl(P_k\,\Pi_l\bigr)^{-\theta\lambda/(1+\theta\lambda)}.
  \label{eq:tkl_reduced}
\end{equation}
The same algebra gives the equilibrium link traffic:
\begin{equation}
  \xi_{kl}
  \;=\;
  \tbar_{kl}^{-\theta/(1+\theta\lambda)}\cdot
  n_{kl}^{\theta\lambda/(1+\theta\lambda)}\cdot
  \bigl(P_k\,\Pi_l\bigr)^{-\theta/(1+\theta\lambda)}.
  \label{eq:xi_reduced}
\end{equation}

\paragraph{Comparative statics.}
Equations \eqref{eq:tkl_reduced}--\eqref{eq:xi_reduced} deliver the
partial (i.e.\ accessibility-held-fixed) elasticities
\begin{equation}
  \frac{\partial\ln\tkl}{\partial\ln n_{kl}}
  \;=\;-\,\frac{\lambda}{1+\theta\lambda}\;<\;0
  \qquad\text{vs.}\qquad
  \frac{\partial\ln\xi_{kl}}{\partial\ln n_{kl}}
  \;=\;\frac{\theta\lambda}{1+\theta\lambda}\;>\;0.
  \label{eq:elasticities}
\end{equation}
\emph{Adding a directed lane reduces the directed travel time on that
link, but because travellers re-route toward the now-faster edge it
also increases the directed traffic.  The two effects jointly
determine the equilibrium adjustment.}  The absolute ratio of the
two elasticities equals $\theta$: the elasticity of directed
traffic with respect to directed lanes is $\theta$ times larger
(in absolute value) than the elasticity of directed travel time
with respect to directed lanes.  In the limit of very strong
congestion ($\theta\lambda\to\infty$), $\partial\ln\xi_{kl}/
\partial\ln n_{kl}\to 1$, matching the ``fundamental law of road
congestion'' of \cite{DurantonTurner2011}.

\paragraph{Tidal-lane asymmetry channel.}
The tidal-lane policy we consider moves directed lanes from the
off-peak to the peak direction on a targeted set of links
$\E^\star\subseteq\E$ while leaving total corridor capacity
$n_{kl}+n_{lk}$ unchanged.  Because each direction enters the
congestion technology \eqref{eq:cong_primitive} separately, the
policy strictly reduces $\tkl$ and strictly increases $t_{lk}$ on
treated corridors, leading to asymmetric changes in bilateral
commuting costs $\{\tijk\}$, bilateral flows $\{L_{ij}\}$, and
residence--workplace allocations $(L^R,L^F)$.

\subsection{Equilibrium}\label{ssec:eq}

\paragraph{Reduced equilibrium system in population shares.}
Let $l_i^R\equiv L_i^R/\Lbar$ and $l_j^F\equiv L_j^F/\Lbar$ denote
residential and workplace population shares.  Combining the
gravity equation \eqref{eq:gravity}, the fundamentals
\eqref{eq:fundamentals}, the wage equation \eqref{eq:wage}, and the
market-clearing conditions \eqref{eq:marketclearing}, and using
$L_i^R=\Lbar\,l_i^R$, $L_j^F=\Lbar\,l_j^F$ delivers the pair of
equilibrium conditions
\begin{align}
  (l_i^R)^{1-\beta\theta}
  &=\chi\sum_{j\in\N}\tijk^{-\theta}\,\ubar_i^{\theta}\,\Abar_j^{\theta}\,
       (l_j^F)^{\alpha\theta},
  \label{eq:eq_R}\\[2pt]
  (l_j^F)^{1-\alpha\theta}
  &=\chi\sum_{i\in\N}\tijk^{-\theta}\,\ubar_i^{\theta}\,\Abar_j^{\theta}\,
       (l_i^R)^{\beta\theta},
  \label{eq:eq_F}
\end{align}
where the scalar
\begin{equation}
  \chi\;\equiv\;\frac{\Lbar^{\theta(\alpha+\beta)}}{\Wbar^{\theta}}
  \label{eq:chi_def}
\end{equation}
is pinned down by a normalization of the fundamentals.  The
step-by-step derivation of \eqref{eq:eq_R}--\eqref{eq:chi_def} is in
Appendix~\ref{app:eqshares}.

\paragraph{Definition of equilibrium.}

\begin{definition}[Equilibrium]\label{def:eq}
Given fundamentals $\{(\ubar_i,\Abar_j)\}$, directed link primitives
$\{(\ell_{kl},v^{0}_{kl},n_{kl})\}_{(k,l)\in\E}$, and parameters
$(\theta,\alpha,\beta,\lambda)$, an \emph{equilibrium} is a tuple
\[
  \Bigl(\{L_i^R,L_j^F\}_{i,j\in\N},\ \{L_{ij}\}_{(i,j)\in\N^2},\
  \{\tijk\}_{(i,j)\in\N^2},\ \{\tkl\}_{(k,l)\in\E},\
  \{\xi_{kl}\}_{(k,l)\in\E}\Bigr)
\]
such that
\begin{enumerate}[label=(\roman*),leftmargin=2em,itemsep=0.2em]
\item given $\{\tkl\}$, the bilateral cost index $\{\tijk\}$ satisfies
  the series representation \eqref{eq:tau_series};
\item given $\{\tijk\}$, bilateral flows $\{L_{ij}\}$ satisfy the
  gravity equation \eqref{eq:gravity};
\item marginal populations $\{L_i^R,L_j^F\}$ satisfy
  \eqref{eq:marketclearing} and the share system
  \eqref{eq:eq_R}--\eqref{eq:eq_F};
\item directed traffic $\{\xi_{kl}\}$ is generated by
  \eqref{eq:xi_def} and equivalently satisfies the link-traffic gravity
  \eqref{eq:xi_gravity};
\item directed travel times $\{\tkl\}$ are consistent with the
  congestion technology \eqref{eq:cong_primitive}.
\end{enumerate}
\end{definition}

\paragraph{Existence and uniqueness.}
Conditions (i)--(iii) of Definition~\ref{def:eq} recover the
Ahlfeldt-type quantitative spatial equilibrium with directional
commuting costs; conditions (iv)--(v) introduce endogenous congestion
and follow \cite{AllenArkolakis2022}.  Under $\lambda\ge 0$ and
Assumption~\ref{ass:specrad}, existence follows from the
Brouwer-based fixed-point argument in Proposition~1 of
\cite{AllenArkolakis2022}.

Proposition~1(b) of \cite{AllenArkolakis2022} supplies a
\emph{sufficient} condition for uniqueness in the urban variant:
\begin{equation}
  \alpha\;\le\;\tfrac{1}{2}\!\left(\tfrac{1}{\theta}-\lambda\right),
  \qquad
  \beta\;\le\;\tfrac{1}{2}\!\left(\tfrac{1}{\theta}-\lambda\right),
  \label{eq:AAunique}
\end{equation}
which must hold jointly.  Since we impose $\alpha>0$ and $\beta<0$
(Eq.~\eqref{eq:sign}), the right-hand inequality is implied by the
sign restriction on $\beta$, and the left-hand inequality on
$\alpha$ is the binding one.  \eqref{eq:AAunique} is a purely
algebraic inequality in $(\alpha,\beta,\theta,\lambda)$; whether it
holds at a given grid point is decided by substitution, not by
numerical verification.  At the baseline $(\theta,\alpha)=(6.83,0.06)$
the binding inequality requires $\lambda\le 1/\theta-2\alpha
\approx 0.026$, which is below every point of the robustness grid
$\lambda\in\{0.1,0.3,0.6,1.0\}$; the sufficient condition therefore
fails throughout the robustness exercise.  At $\alpha=0.08$ in the externality-grid
robustness, the required bound is $\lambda\le -0.014$, which no
non-negative $\lambda$ satisfies; the sufficient condition therefore
fails throughout the admissible $\lambda\ge 0$ range at that grid
point.

Failure of \eqref{eq:AAunique} is not a statement about
non-uniqueness, only about the reach of the Allen-Arkolakis sufficient
bound.  On grid points where \eqref{eq:AAunique} fails, uniqueness of
the numerical solution is supported by three checks, reported in
Appendix~\ref{app:uniqueness_numeric}: (i) the solver converges to
the same fixed point from multiple randomised initialisations;
(ii) the spectral radius of the Jacobian of the equilibrium map at
the fixed point is strictly less than one; (iii) the equilibrium
varies continuously with the parameters across the grid.  In every
specification we examine the numerical equilibrium is unique and
stable under iteration.  No formal claim of uniqueness beyond the
range where \eqref{eq:AAunique} holds is made.

\begin{remark}[Degenerate cases]\label{rem:degenerate}
If $\lambda=0$, conditions (iv)--(v) are vacuous, $\tkl=\tbar_{kl}$,
and the equilibrium reduces to a pure spatial model with directional
but exogenous commuting costs.  If additionally the network is
undirected (so $\tbar_{kl}=\tbar_{lk}$), $\tijk=\tau_{ji}$ and the
model collapses to a textbook Ahlfeldt et al.\ spatial equilibrium.
\end{remark}

\subsection{Welfare}\label{ssec:welfare}

This section derives the welfare representation that we use both as a
positive object (an equilibrium summary statistic) and as a normative
object (the planner's objective in counterfactual evaluations).  All
derivations follow from primitives already introduced; detailed
proofs are collected in Appendix~\ref{app:welfare2}.

\paragraph{Ex-ante expected utility.}
Recall from \S\ref{ssec:ind} that, conditional on
$(\{u_i\},\{w_j\},\{\tkl\})$, the realised utility of a worker with
choice $(i,j,r)$ is $V_{ij,r}$ in \eqref{eq:utility}.  Each worker
chooses $(i,j,r)$ to maximise $V$.  We define \emph{social welfare}
$\mathbb W$ as the ex-ante expected value of the maximised utility,
taken over the i.i.d.\ Fr\'echet shocks:
\begin{equation}
  \mathbb W
  \;\equiv\;
  \mathbb E\!\left[\,\max_{(i,j,r)\in\N^{2}\times\R_{ij}}V_{ij,r}\,\right].
  \label{eq:welfare_def}
\end{equation}
Because $\varepsilon_{ij,r}$ are i.i.d.\ Fr\'echet$(\theta)$ and the
deterministic component is $\omega_{ij,r}=u_iw_j/\prod_m
t_{r_{m-1},r_m}$, the maximum $M\equiv\max_{(i,j,r)}V_{ij,r}$ is
itself Fr\'echet with shape $\theta$ and scale $S^{1/\theta}=\Wbar$
defined in \eqref{eq:S_def}--\eqref{eq:W_def}.  Therefore
\begin{equation}
  \boxed{\;\mathbb W
  \;=\;\Gamma\!\left(1-\tfrac{1}{\theta}\right)\,\Wbar
  \;=\;\Gamma\!\left(1-\tfrac{1}{\theta}\right)
       \left[\sum_{(i,j)\in\N^{2}}
             \tijk^{-\theta}\,u_i^{\theta}\,w_j^{\theta}\right]^{1/\theta}.\;}
  \label{eq:welfare_main}
\end{equation}
Equation \eqref{eq:welfare_main} establishes that welfare is, up to the
dispersion-only constant $\Gamma(1-1/\theta)$, equal to the aggregate
commuter-market-access scale $\Wbar$.  The dimensional factor
$\Gamma(1-1/\theta)$ is required for $\theta>1$; it absorbs the
location of the Fr\'echet mean and is invariant to model primitives.

\paragraph{Market-access representation of welfare.}
Combining \eqref{eq:welfare_main} with the resident and firm
market-access indices \eqref{eq:Pi_def}--\eqref{eq:Pj_def} delivers
two equivalent expressions for $\Wbar^{\theta}$:
\begin{equation}
  \Wbar^{\theta}
  \;=\;\sum_{i\in\N}u_i^{\theta}\,\Pi_i^{-\theta}
  \;=\;\sum_{j\in\N}w_j^{\theta}\,P_j^{-\theta}.
  \label{eq:welfare_MA}
\end{equation}
The first equality says aggregate welfare is the residence-weighted
sum of resident-side amenities and accessibility; the second says it
is the workplace-weighted sum of firm-side wages and accessibility.
Equality across the two follows from the commuting market-clearing
condition \eqref{eq:marketclearing} and provides a useful internal
consistency check in numerical implementations.

\paragraph{Per-capita and per-pair welfare.}
Because every worker draws an identically distributed Fr\'echet shock,
\eqref{eq:welfare_main} is also the \emph{per-capita} ex-ante welfare:
$\mathbb W=\mathbb E[V^{\max}_{n}]$ for any individual worker $n$.
Conditional on having chosen residence--workplace pair $(i,j)$, the
expected realised utility is
\begin{equation}
  \mathbb E\bigl[V_{ij}\,\big|\,\text{chose }(i,j)\bigr]
  \;=\;\mathbb W
  \;=\;\Gamma\!\left(1-\tfrac{1}{\theta}\right)\Wbar,
  \label{eq:welfare_constant_across_pairs}
\end{equation}
i.e.\ \emph{indifference across chosen pairs}: every commuting pair
delivers the same expected utility ex post, which is the standard
Fr\'echet-spatial-equilibrium result and underlies the use of $\Wbar$
as a single welfare scalar.\footnote{Heterogeneity in welfare arises
only across realisations of $\varepsilon_{ij,r}$; the population
averages out the idiosyncratic component.}

\paragraph{Welfare decomposition.}
Differentiating $\ln\mathbb W=\ln\Gamma(1-1/\theta)+\ln\Wbar$ and
using \eqref{eq:welfare_main} gives the exact log-decomposition
\begin{equation}
  \mathrm d\ln\mathbb W
  \;=\;
  \underbrace{\sum_i\,\Lambda_i^{R}\,\mathrm d\ln u_i}_{\text{amenity channel}}
  \;+\;
  \underbrace{\sum_j\,\Lambda_j^{F}\,\mathrm d\ln w_j}_{\text{productivity channel}}
  \;-\;
  \underbrace{\sum_{(i,j)}\,\Lambda_{ij}\,\mathrm d\ln\tijk}_{\text{access channel}},
  \label{eq:welfare_decomp}
\end{equation}
where the weights are baseline commuting shares,
\begin{equation}
  \Lambda_i^{R}\;\equiv\;\frac{L_i^R}{\Lbar},
  \qquad
  \Lambda_j^{F}\;\equiv\;\frac{L_j^F}{\Lbar},
  \qquad
  \Lambda_{ij}\;\equiv\;\frac{L_{ij}}{\Lbar},
  \label{eq:lambda_weights}
\end{equation}
and $\sum_i\Lambda_i^R=\sum_j\Lambda_j^F=\sum_{(i,j)}\Lambda_{ij}=1$.
Equation \eqref{eq:welfare_decomp} provides a transparent
accounting: the welfare effect of any policy or shock is the
flow-weighted average of its impact on local amenities ($u_i$), local
productivity ($w_j$), and bilateral commuting costs ($\tijk$).  In
the presence of endogenous spillovers \eqref{eq:fundamentals}, the
amenity and productivity channels themselves respond to changes in
$(L^R,L^F)$ at elasticities $\beta$ and $\alpha$, which generates the
indirect welfare effect that we isolate in
\S\ref{ssec:welfare_hat} via exact-hat algebra.

\begin{remark}[The directional content of welfare]\label{rem:directional_welfare}
A central feature of our setting is that $\tijk\ne\tau_{ji}$ and
$\tkl\ne t_{lk}$.  Within the welfare decomposition
\eqref{eq:welfare_decomp}, this directionality enters in two distinct
ways.  First, the access channel weights $\Lambda_{ij}$ are themselves
asymmetric across $(i,j)$: AM-peak commuting flows are
predominantly suburb-to-CBD, so the marginal welfare value of a
reduction in inbound peak-direction travel time is, in general,
strictly larger than the marginal welfare cost of an equivalent
increase in outbound off-peak travel time.  Second, a tidal-lane
reallocation that shifts capacity $n_{kl}$ from one direction to its
counter-direction enters \eqref{eq:welfare_decomp} through
\emph{both} $\mathrm d\ln\tau_{kl}$ and $\mathrm d\ln\tau_{lk}$ but
with opposite signs and unequal weights.  The net welfare effect of
a tidal-lane policy is therefore the difference of two flow-weighted
accessibility terms and is generically non-zero, even when total
lane-kilometres are conserved.  Quantifying this difference is the
central counterfactual exercise of the paper.
\end{remark}

\paragraph{Hat algebra for welfare.}\label{ssec:welfare_hat}
For counterfactual policy analysis it is convenient to express
changes in $\mathbb W$ as a function of \emph{relative} changes
$\hat x\equiv x'/x$ from the baseline equilibrium.  Define
\begin{equation}
  \hat{\mathbb W}\;\equiv\;\frac{\mathbb W'}{\mathbb W}
  \;=\;\frac{\Wbar'}{\Wbar}\;\equiv\;\hat{\Wbar},
  \label{eq:Whatdef}
\end{equation}
where the dispersion constant $\Gamma(1-1/\theta)$ cancels.  Using
\eqref{eq:welfare_main} and the gravity equation
\eqref{eq:gravity} to eliminate $u_i^{\theta}w_j^{\theta}\tijk^{-\theta}$
in favour of baseline commuting shares
$\pi_{ij}^{\obs}\equiv L_{ij}^{\obs}/\Lbar$, one obtains the
\emph{exact-hat welfare formula}
\begin{equation}
  \boxed{\;\hat{\mathbb W}^{\theta}
  \;=\;\hat{\Wbar}^{\theta}
  \;=\;\sum_{(i,j)\in\N^{2}}
        \pi_{ij}^{\obs}\,
        \hat{u}_i^{\theta}\,\hat{w}_j^{\theta}\,\hat{\tau}_{ij}^{-\theta}.\;}
  \label{eq:welfare_hat}
\end{equation}
With endogenous spillovers \eqref{eq:fundamentals} the relative
fundamentals are pinned down by the share changes,
\begin{equation}
  \hat u_i\;=\;\hat l_i^{R,\beta},
  \qquad
  \hat w_j\;=\;\hat l_j^{F,\alpha},
  \label{eq:hat_uw}
\end{equation}
so that \eqref{eq:welfare_hat} can be evaluated using only baseline
commuting shares $\{\pi_{ij}^{\obs}\}$, the counterfactual link
times $\{\hat{\tau}_{ij}\}$ (which obtain from
\S\ref{ssec:hat}--Appendix~\ref{app:hat}), and the share changes
$\{\hat l_i^R,\hat l_j^F\}$.  Critically, no knowledge of the
\emph{levels} of $(\ubar_i,\Abar_j)$ is required for welfare changes;
this is the exact-hat property exploited by the
Allen--Arkolakis class.

\begin{proposition}[Welfare bounds from market-access changes]
\label{prop:welfare_bounds}
Define the residence-side and workplace-side hat market-access
indices by
$\hat{\Pi}_i^{-\theta}\equiv\sum_{j}\pi_{ij|i}^{\obs}\hat{w}_j^{\theta}\hat{\tau}_{ij}^{-\theta}$
and
$\hat{P}_j^{-\theta}\equiv\sum_{i}\pi_{ij|j}^{\obs}\hat{u}_i^{\theta}\hat{\tau}_{ij}^{-\theta}$,
where $\pi_{ij|i}^{\obs}\equiv L_{ij}/L_i^R$ and
$\pi_{ij|j}^{\obs}\equiv L_{ij}/L_j^F$ are conditional commuting
shares.  Then
\begin{equation}
  \hat{\mathbb W}^{\theta}
  \;=\;\sum_{i}\Lambda_i^R\,\hat u_i^{\theta}\,\hat\Pi_i^{-\theta}
  \;=\;\sum_{j}\Lambda_j^F\,\hat w_j^{\theta}\,\hat P_j^{-\theta}.
  \label{eq:welfare_hat_MA}
\end{equation}
\end{proposition}
\noindent The proof, again, is in
Appendix~\ref{app:welfare2}.  The two equivalent representations
\eqref{eq:welfare_hat_MA} are useful in numerical work: they provide
two independent computations of $\hat{\mathbb W}$ from the same
counterfactual block, and disagreement between them flags a
mis-specification or numerical bug.

\begin{remark}[Equivalent variation in time units]
Because commuting cost enters utility multiplicatively as
$\prod_m t^{-1}$, and utility itself is dimensionless once
normalised by $\Gamma(1-1/\theta)$, the welfare change
$\hat{\mathbb W}$ has a natural interpretation as a
``time-equivalent variation''.  Specifically, a uniform proportional
change $\hat{\tau}_{ij}=\kappa$ for all $(i,j)$ (holding $u_i,w_j$
fixed) yields $\hat{\mathbb W}=\kappa^{-1}$.  A counterfactual that
generates $\hat{\mathbb W}=1.02$ is therefore equivalent in welfare
terms to a uniform $1.96\%$ reduction in every bilateral commuting
time.\footnote{$\kappa^{-1}=1.02\Rightarrow \kappa\approx 0.9804$,
i.e.\ a $1.96\%$ decline in $\tijk$.}  This metric is reported
alongside dollarised valuations in
Section~\ref{ssec:cf}--Appendix~\ref{app:hat}.
\end{remark}

\begin{remark}[Welfare under alternative normalisations]
The welfare formula \eqref{eq:welfare_hat} is invariant to the
normalisations $\prod_i\ubar_i=1$ and $\prod_j\Abar_j=1$ used in the
inversion of fundamentals (\S\ref{ssec:invert}), because both
$\mathbb W'$ and $\mathbb W$ scale by the same constant under any
change of normalisation.  This invariance is what allows us to read
$\hat{\mathbb W}$ off the counterfactual block without taking a
stance on the level of $(\ubar_i,\Abar_j)$.
\end{remark}

%====================================================================
\section{Estimation}\label{sec:est}
%====================================================================

\subsection{Data-to-Model Mapping}\label{ssec:data}

We implement the model on a square 3\,km grid of Beijing and use the
AM peak (morning rush) as the main period.  A parallel pipeline for
the PM peak is maintained for robustness, and hexagonal and Voronoi
grids are produced as alternative discretizations.  All empirical
objects below are constructed identically across periods and grid
systems.

\paragraph{Grid-level network $(\N,\E)$.}
Each spatial cell is a node $i\in\N$ identified by a unique
\texttt{grid\_id} and endowed with a geographic centroid.
A directed link $(k,l)$ is included in $\E$ if and only if the
underlying centerline network supports at least one matched
centerline piece whose travel runs from (a location inside) grid $k$
to (a location inside) grid $l$.  The matching pipeline produces, for
each raw centerline piece, (i) its direction, (ii) the sequence of
grids it traverses, (iii) its length inside each traversed grid, and
(iv) its average AM-peak and nighttime speeds; directed grid links
aggregate over all pieces that cross the relevant grid boundary.

\paragraph{Directed link primitives $(\ell_{kl},v^{0}_{kl},n_{kl})$.}
For each directed grid link $(k,l)\in\E$ we construct:
\begin{itemize}[leftmargin=2em,itemsep=0.2em]
  \item \textbf{Grid distance} $\ell_{kl}$: centroid-to-centroid
        distance between adjacent grids $k$ and $l$
        (directional after accounting for one-way road segments);
  \item \textbf{Observed (congestion-inclusive) speed}
        $v_{kl}^{\obs}$: length-weighted harmonic mean of the AM-peak
        speeds of all centerline pieces that cross from $k$ to $l$;
  \item \textbf{Free-flow speed} $v_{kl}^{0}$: length-weighted
        harmonic mean of centerline speeds observed in the
        22:00--05:00 window, aggregated first to centerlines and then
        to the grid link; this is our empirical proxy for the
        non-congested baseline;
  \item \textbf{Directed lane count} $n_{kl}$: length-weighted
        average of approximate per-segment lane counts derived from
        the matched centerlines' \texttt{roadtype} attribute via a
        fixed roadtype-to-lanes lookup table.
\end{itemize}
Observed and free-flow directed travel times follow:
\begin{equation}\label{eq:data_tkl}
  \tkl^{\obs}\;=\;\frac{\ell_{kl}}{v_{kl}^{\obs}},
  \qquad
  \tbar_{kl}\;=\;\frac{\ell_{kl}}{v_{kl}^{0}},
\end{equation}
measured in minutes.  Directionality is preserved throughout: all of
$\ell_{kl}$, $v_{kl}^{\obs}$, $v_{kl}^{0}$, and $n_{kl}$ may differ
from their $(l,k)$ counterparts.

\paragraph{Bilateral commuting cost $\tijk$.}
Given $\{\tkl^{\obs}\}$, we compute the network cost index $\tijk$ in
two complementary ways:
\begin{enumerate}[label=(\alph*),leftmargin=2em,itemsep=0.2em]
  \item \textbf{Shortest-path time}: $\hat\tau_{ij}^{\mathrm{SP}}
        =\min_{p\in\mathcal P_{ij}}\sum_{(k,l)\in p}\tkl^{\obs}$,
        where $\mathcal P_{ij}$ is the set of simple directed paths
        from $i$ to $j$ on $(\N,\E)$;
  \item \textbf{Series cost index}: $\hat\tau_{ij}^{\mathrm{SERIES}}
        =\bigl([(I-A^{\obs})^{-1}]_{ij}\bigr)^{-1/\theta}$, using
        $A^{\obs}_{kl}=\tkl^{\obs,-\theta}$.
\end{enumerate}
Method~(a) is computationally cheap and enters reduced-form
diagnostics; method~(b) is the structural cost index and is used in
the estimation and counterfactual routines.  Reachability is enforced
by restricting attention to the largest weakly connected component
of the directed grid graph.

\paragraph{Commuting flows and marginal populations.}
Using the Baidu Map commuting matrix, we assign each observed
home coordinate and work coordinate to a grid cell via
point-in-polygon matching and aggregate by $(i,j)$:
\begin{equation}\label{eq:data_Lij}
  L_{ij}^{\obs}
  \;\equiv\;\sum_{\text{commuters }(i\to j)}\mathrm{pop},
\end{equation}
where $\mathrm{pop}$ is the number of commuters associated with each
home--work pair.  Residential and workplace marginal totals are
\begin{equation}\label{eq:data_margins}
  L_i^{R,\obs}=\sum_{j\in\N}L_{ij}^{\obs},\quad
  L_j^{F,\obs}=\sum_{i\in\N}L_{ij}^{\obs},\quad
  \Lbar=\sum_i L_i^{R,\obs}=\sum_j L_j^{F,\obs}.
\end{equation}
When available we retain commute-mode counts (walk, bike, subway,
bus, car) for descriptive purposes.

\paragraph{Empirical link traffic $\hat\xi_{kl}$.}
Because our data do not contain directly observed flows on directed
links, we \emph{construct} $\hat\xi_{kl}$ by probabilistic assignment
of observed OD flows onto the directed grid graph.  Following the
Fr\'echet route-choice model, the share of commuters from $(i,j)$
passing through $(k,l)$ is $\pi^{kl}_{ij}$ in \eqref{eq:piijkl}, so
\begin{equation}\label{eq:hatxi}
  \hat\xi_{kl}
  \;\equiv\;
  \sum_{i,j\in\N}L_{ij}^{\obs}\,\pi^{kl}_{ij}(\theta,\{\tkl^{\obs}\}).
\end{equation}
By Proposition~\ref{prop:linkgravity}, \eqref{eq:hatxi} is model
consistent: if the observed OD matrix were an equilibrium, the
implied flows would satisfy \eqref{eq:xi_gravity}.  In our baseline
we use \eqref{eq:hatxi} with the externally calibrated $\theta$; an
all-or-nothing shortest-path assignment is retained as a
robustness check.

\paragraph{Summary table.}

\begin{table}[h!]
\centering
\caption{Model variables and empirical construction.}\label{tab:data_mapping}
\small
\begin{tabular}{p{3.6cm}lp{7.8cm}}
\toprule
Variable & Symbol & Construction \\
\midrule
Location (node) & $i\in\N$ & Grid cell centroid (square 3\,km baseline) \\
Directed edge & $(k,l)\in\E$ & Adjacency of grid cells supported by matched centerline data \\
Grid distance & $\ell_{kl}$ & Centroid-to-centroid distance between $k,l$ \\
Observed directed speed & $v_{kl}^{\obs}$ & Length-weighted harmonic mean of AM speeds of centerline pieces crossing $(k,l)$ \\
Free-flow directed speed & $v_{kl}^{0}$ & Same aggregator applied to 22:00--05:00 pieces \\
Directed lane count & $n_{kl}$ & Length-weighted lane count from roadtype lookup \\
Observed travel time & $\tkl^{\obs}$ & $\ell_{kl}/v_{kl}^{\obs}$ \\
Free-flow travel time & $\tbar_{kl}$ & $\ell_{kl}/v_{kl}^{0}$ \\
Bilateral cost index & $\tijk$ & $[(I-A^{\obs})^{-1}]_{ij}^{-1/\theta}$; shortest path as diagnostic \\
Commuting flow & $L_{ij}^{\obs}$ & Baidu home--work matrix aggregated to grids \\
Residential / workplace totals & $L_i^{R,\obs},L_j^{F,\obs}$ & Row / column sums of $L_{ij}^{\obs}$ \\
Link traffic & $\hat\xi_{kl}$ & Probabilistic OD-to-link assignment via $\pi^{kl}_{ij}$ \\
Baseline fundamentals & $(\ubar_i,\Abar_j)$ & Inverted from \eqref{eq:eq_R}--\eqref{eq:eq_F} given other inputs \\
\bottomrule
\end{tabular}
\end{table}

\paragraph{Reachability and sample.}
Before estimation we restrict attention to the set of grids that lie
in the largest weakly connected component of $(\N,\E)$ and to OD
pairs whose origin and destination both lie in this component.  In
the square 3\,km grid, the raw Baidu commuting matrix contains
$442{,}796$ distinct origin--destination pairs and $7.96$ million
commuters; the reachability filter retains $374{,}850$ pairs
($84.7\%$) covering $7.59$ million commuters ($95.3\%$).  This step
eliminates peripheral, low-volume pairs that would otherwise
contaminate the cost-index and market-access computations.

\subsection{Calibrated and Externally Disciplined Parameters}\label{ssec:calib}

The baseline implementation calibrates $(\theta,\alpha,\beta)$
externally and estimates $\lambda$ from the data; the fundamentals
$(\ubar_i,\Abar_j)$ are then inverted.  A gravity-based internal
estimate of $\theta$ is produced as a robustness check
(Appendix~\ref{app:theta}).

\paragraph{$\theta$ (Fr\'echet dispersion).}
In the main text we set $\theta=6.83$, the value estimated by
\cite{Ahlfeldt2015} in a closely related quantitative spatial
equilibrium with endogenous commuting.  This value lies inside the
range routinely used in the quantitative urban literature
($\theta\in[4,9]$) and has the additional advantage of making our
numbers directly comparable to \cite{AllenArkolakis2022} and
\cite{Ahlfeldt2015}.

\paragraph{$\alpha,\beta$ (externality elasticities).}
Under the AA-style one-parameter specification of
\S\ref{ssec:env}--Remark~\ref{rem:AA_vs_Ahlfeldt}, we set
\begin{equation}
  \alpha=0.06,
  \qquad
  \beta=-0.30,
  \label{eq:alphabeta_bench}
\end{equation}
consistent with the sign restrictions $\alpha>0$, $\beta<0$.  The
value $\alpha=0.06$ is the central estimate of the production
agglomeration elasticity in the meta-style survey of
\cite{CombesGobillon2015}, falls within the $[0.03,0.10]$ range
documented across developed and developing cities
(\cite{CombesDurantonGobillon2008,CombesGobillon2015}), and is the
benchmark used by \cite{AllenArkolakis2014,AllenArkolakis2022} in
their quantitative evaluations.  The value $\beta=-0.30$ is the
residential-externality elasticity used by
\cite{AllenArkolakis2014}; it is consistent with the implied
residential spillover magnitudes in
\cite{MonteReddingRossiHansberg2018} and
\cite{HeblichReddingSturm2020}.%
\footnote{%
We deliberately do \emph{not} borrow $(\alpha,\beta)$ from
\cite{Ahlfeldt2015}, whose reduced-form composite parameters are
constructed under an explicit two-force (agglomeration plus
congestion) specification and therefore have a different
interpretation than the one-parameter composite used here (see
Remark~\ref{rem:AA_vs_Ahlfeldt}).  Using AA-class benchmarks keeps
the interpretation of the externality elasticities internally
consistent with the rest of the model.}
These magnitudes generate meaningful but non-dominant local
spillovers.  In Section~\ref{ssec:cf} we report sensitivity to
$(\alpha,\beta)$ over the grids $\alpha\in\{0.04,0.06,0.08\}$ and
$\beta\in\{-0.50,-0.30,-0.20\}$, which bracket the AA-class estimates
in the literature.

\paragraph{Verification of equilibrium existence conditions.}
The sufficient conditions for equilibrium existence
(Definition~\ref{def:eq} and the accompanying fixed-point
argument) require $\theta\alpha<1$, $\theta\beta<1$, $\lambda\ge 0$,
and the spectral radius condition $\spr(A)<1$
(Assumption~\ref{ass:specrad}).  At the benchmark calibration
\begin{equation*}
  \theta\alpha=6.83\times 0.06\approx 0.410<1,
  \qquad
  \theta\beta=6.83\times(-0.30)\approx -2.049<1,
\end{equation*}
so conditions (i) and (ii) are satisfied with comfortable slack; the
spectral-radius condition on $A$ is checked numerically on the
solution grid (it holds in every specification we examine).  Both
robustness grids $\alpha\in\{0.04,0.06,0.08\}$ and
$\beta\in\{-0.50,-0.30,-0.20\}$ also satisfy these conditions.

\paragraph{$\lambda$ (congestion elasticity).}
Identified from an instrumental-variable regression of directed
speed on directed traffic per lane; details in
Section~\ref{ssec:lambdahat}.  OLS on the same equation is
mechanically circular (the right-hand regressor is built from the
dependent variable through the link cost
$\tkl^{\obs}=\ell_{kl}/v_{kl}^{\obs}$) and is reported only for
comparison.  The IV design is in Appendix~\ref{app:iv}.

\paragraph{Normalization.}
The scalar $\chi$ in \eqref{eq:chi_def} is not separately identified
from the level of $(\ubar_i,\Abar_j)$.  We normalize by imposing
$\prod_{i}\ubar_i=1$ and $\prod_{j}\Abar_j=1$ in the inversion
(Section~\ref{ssec:invert}); equivalently, $\Wbar^{\theta}$ is set
such that $\sum_i l_i^R=\sum_j l_j^F=1$ in the baseline equilibrium.

\subsection{Congestion Elasticity $\lambda$: Reduced-Form Estimation}\label{ssec:lambdahat}

We estimate $\lambda$ from a reduced-form speed--flow regression
whose specification is derived from the congestion technology
\eqref{eq:cong_primitive}.

\paragraph{Mapping from the congestion equation to a speed--flow
regression.}
Substituting $\tkl=\ell_{kl}/v_{kl}$ and $\tbar_{kl}=\ell_{kl}/v^{0}_{kl}$
into \eqref{eq:cong_primitive} and taking logs gives the identity
\begin{equation}\label{eq:speed_identity}
  \ln v_{kl}
  \;=\;
  \ln v^{0}_{kl}
  -\lambda\,\ln\!\left(\frac{\xi_{kl}}{n_{kl}}\right).
\end{equation}
Relaxing the constraint that the intercept equals $\ln v^{0}_{kl}$
exactly (so that unobserved, link-specific capacity effects are
absorbed), we obtain the estimable equation
\begin{equation}\label{eq:lambda_regression}
  \boxed{\;
  \ln v_{kl}^{\obs}
  \;=\;
  m_0
  +\delta_1\,\ln\!\left(\frac{\hat\xi_{kl}}{n_{kl}}\right)
  +\eta_{\mathrm{bin}(kl)}
  +\varepsilon_{kl},
  \qquad
  \hat\lambda=-\hat\delta_1.
  \;}
\end{equation}
Here $m_0$ captures an average free-flow log-speed; $\eta_{\mathrm{bin}(kl)}$
is a set of fixed effects that absorb systematic variation in
baseline capacity across link types (e.g.\ lane-count bin,
roadtype-class, or CBD distance bin); $\hat\xi_{kl}$ is constructed
by the probabilistic OD-to-link assignment \eqref{eq:hatxi}; and
$\varepsilon_{kl}$ is a mean-zero disturbance capturing measurement
error in speeds, lane counts, or the flow construction.

\paragraph{Relation to the handwritten-note speed equation.}
The handwritten notes specify
$\ln v_{kl}=\ln m_0+\delta_1\ln(\xi_{kl}/n_{kl})+\ln\varepsilon_{kl}$
and, under a generic iceberg micro-foundation
$\tkl=(\ell_{kl}/v_{kl})^{\delta_0}$ with
$v_{kl}=v^{0}_{kl}(n_{kl}/\xi_{kl})^{\lambda/\delta_0}$, deliver the
mapping $\lambda=-\delta_0\cdot\delta_1$.  The physical-time
convention $\delta_0=1$ adopted in Section~\ref{ssec:ind}
(Remark~\ref{rem:iceberg}) collapses this to
$\hat\lambda=-\hat\delta_1$, which coincides with the boxed
estimating equation \eqref{eq:lambda_regression}.  Had the
\cite{AllenArkolakis2022} convention $\delta_0=1/\theta$ been adopted
instead, the mapping would read $\hat\lambda^{\mathrm{AA}}
=-\hat\delta_1/\theta$, smaller by a factor of $\theta$ but referring
to a renormalised iceberg-cost elasticity rather than to a
physical-time elasticity.  The choice between the two is dimensional;
it does not affect equilibrium predictions, and it does affect the
numerical value reported for $\hat\lambda$ by exactly the factor
$\theta$.

\paragraph{Construction of the regressor $\hat\xi_{kl}$.}
The default assignment \eqref{eq:hatxi} builds $\hat\xi_{kl}$ from
route probabilities $\pi^{kl}_{ij}$ whose underlying bilateral cost
index is a function of $\{\tkl^{\obs}\}$, and
$\tkl^{\obs}=\ell_{kl}/v_{kl}^{\obs}$ is itself a function of the
dependent variable $v_{kl}^{\obs}$.  Regressing $\ln v_{kl}^{\obs}$
on $\ln(\hat\xi_{kl}/n_{kl})$ under this construction is mechanically
circular: the regressor absorbs the same realisation of $v_{kl}^{\obs}$
that generates the disturbance $\varepsilon_{kl}$, so the estimand no
longer corresponds to the congestion elasticity defined by
\eqref{eq:cong_primitive}.  The resulting bias on $\hat\lambda$ has
no sign restriction that flips with the direction of the externality
and is, in particular, not bounded by it.

We therefore construct $\hat\xi_{kl}$ in three ways and report all
three.
\begin{enumerate}[label=(\alph*),leftmargin=2em,itemsep=0.15em]
  \item \emph{Free-flow (baseline OLS).}  Replace $\{\tkl^{\obs}\}$
        by $\{\tbar_{kl}=\ell_{kl}/v^{0}_{kl}\}$ in the route-probability
        calculation, so that $\pi^{kl}_{ij}$ depends only on nighttime
        free-flow link costs and network geometry.  The resulting
        $\hat\xi_{kl}^{\text{ff}}$ is measurable with respect to
        pre-congestion primitives and is orthogonal to
        contemporaneous morning-peak speed shocks by construction.
  \item \emph{Leave-one-out.}  Recompute $\pi^{kl}_{ij}$ with link
        $(k,l)$ excluded from the cost index, so that
        $\hat\xi_{kl}^{\text{loo}}$ does not use $v_{kl}^{\obs}$ on
        the link being regressed; contemporaneous costs on all other
        links are preserved.
  \item \emph{Naive (for comparison only).}  The original
        construction based on $\{\tkl^{\obs}\}$, reported to
        document the magnitude of the circularity bias.
\end{enumerate}
The baseline specification uses $\theta=6.83$ throughout.  The
fixed effects absorb systematic variation in free-flow capacity
across link types but do not resolve the circularity in (c);
resolving it requires construction (a), (b), or the IV of
Appendix~\ref{app:iv}.

\paragraph{Estimation sample and specification choices.}
\begin{enumerate}[label=(\alph*),leftmargin=2em,itemsep=0.15em]
  \item Sample: directed grid links $(k,l)\in\E$ in the largest
        weakly connected component of the AM-peak graph, with at
        least one matched centerline piece and non-zero assigned
        flow;
  \item Weighting: OLS weighted by centerline length share (or
        unweighted; both are reported);
  \item Fixed effects: three lane-count bins (1--2, 3--4, 5+) and
        four roadtype classes (major arterial, minor arterial, local
        road, other), interacted;
  \item Clustering: standard errors clustered at the centerline
        parent (or corridor) level.
\end{enumerate}

\paragraph{Identification assumption and mechanical circularity.}
$\lambda$ is the partial elasticity of travel time with respect to
traffic per lane, holding all other link attributes fixed.  The OLS
identifying condition on \eqref{eq:lambda_regression} reads
\[
  \mathbb E\!\left[\varepsilon_{kl}\,\bigm|\,\ln(\hat\xi_{kl}/n_{kl}),\
  \eta_{\mathrm{bin}(kl)}\right]=0.
\]
When $\hat\xi_{kl}$ is built from $\{\tkl^{\obs}\}$ as in the naive
construction (c) above, this condition fails by construction, not by
omission of controls: $v_{kl}^{\obs}$ enters both sides of the
regression through the dependence of $\hat\xi_{kl}$ on
$\tkl^{\obs}=\ell_{kl}/v_{kl}^{\obs}$.  The resulting $\hat\delta_1$
then captures an accounting relationship between speed and its own
reciprocal rather than the causal elasticity of travel time with
respect to traffic per lane, and the implied bias on $\hat\lambda$ is
not bounded by the sign or magnitude of the externality.

Under constructions (a) or (b) of the previous paragraph, the
regressor no longer embeds $v_{kl}^{\obs}$; what remains is a
conventional simultaneity concern, namely that residual speed
variation may still be correlated with realised traffic per lane
through equilibrium sorting of commuters across routes.  This
residual simultaneity is addressed by the IV design in
Appendix~\ref{app:iv}, which instruments $\ln(\hat\xi_{kl}/n_{kl})$
with link-geography primitives that shift traffic per lane without
affecting contemporaneous speed shocks.  The IV estimator is the
baseline estimator of $\lambda$ in this paper; the OLS variants
under (a)--(c) are reported as comparison rows and as diagnostics.

\paragraph{Baseline estimate.}
The central calibration of $\lambda$ for the counterfactual solver
is the IV estimate $\hat\lambda^{\mathrm{IV}}$
[\textsc{todo: verify point estimate and first-stage $F$ once IV
pipeline finalised}].  Among the OLS rows, the naive construction~(c)
produces $\hat\lambda^{\mathrm{OLS,naive}}$ on the order of
$0.01$--$0.05$ in the square 3\,km grid; because of the mechanical
circularity documented above, the distance of this estimate from the
IV row is reported as a diagnostic of the circularity bias rather
than as an estimand of the structural $\lambda$.  The free-flow
(a) and leave-one-out (b) OLS variants are reported as intermediate
comparison rows.  Counterfactual robustness is conducted at
$\lambda\in\{0.1,0.3,0.6,1.0\}$ in Section~\ref{ssec:cf}.  These values
bracket both the naive-OLS range ($0.01$--$0.05$) and the
Allen--Arkolakis estimate converted to our physical-time convention:
their reported $\lambda^{AA}\approx 0.08$--$0.10$ corresponds to
$\lambda=\lambda^{AA}\cdot\theta\approx 0.55$--$0.68$ in physical-time
units (Section~\ref{rem:iceberg}), which falls near the interior of our
robustness grid.

\subsection{Inversion of Fundamentals $(\ubar_i,\Abar_j)$}\label{ssec:invert}

Given the calibrated parameters $(\theta,\alpha,\beta,\lambda)$ and
the data $(L_i^{R,\obs},L_j^{F,\obs},\{\tkl^{\obs}\})$, we invert
baseline fundamentals $(\ubar_i,\Abar_j)$ such that the observed
spatial allocation is consistent with equilibrium.

\paragraph{Identifying assumption.}
Following the standard inversion logic of quantitative spatial
models, we assume that the observed residential and workplace
populations constitute the baseline equilibrium.

\begin{assumption}[Observed allocation is an equilibrium]
\label{ass:obs_eq}
$(L_i^{R,\obs},L_j^{F,\obs})$ equals $(L_i^R,L_j^F)$ in the
definition of equilibrium (Definition~\ref{def:eq}) under some
baseline fundamentals $(\ubar_i,\Abar_j)$ and under the observed
link primitives $(\ell_{kl},v_{kl}^{\obs},n_{kl})$ with
$\tkl^{\obs}$ given by \eqref{eq:data_tkl}.  The bilateral cost
index $\tijk^{\obs}$ used in the inversion is the series cost
$[(I-A^{\obs})^{-1}]_{ij}^{-1/\theta}$.
\end{assumption}

Under Assumption~\ref{ass:obs_eq}, the equilibrium system
\eqref{eq:eq_R}--\eqref{eq:eq_F} becomes, in the observed data,
\begin{align}
  (l_i^{R,\obs})^{1-\beta\theta}
  &=\chi\sum_{j}(\tijk^{\obs})^{-\theta}\ubar_i^{\theta}\Abar_j^{\theta}
       (l_j^{F,\obs})^{\alpha\theta},
  \label{eq:invert_R}\\
  (l_j^{F,\obs})^{1-\alpha\theta}
  &=\chi\sum_{i}(\tijk^{\obs})^{-\theta}\ubar_i^{\theta}\Abar_j^{\theta}
       (l_i^{R,\obs})^{\beta\theta},
  \label{eq:invert_F}
\end{align}
which is a two-sided nonlinear system in $(\ubar_i,\Abar_j)$ indexed
by $i,j\in\N$.

\paragraph{Algorithm (iterative fixed-point inversion).}

\begin{enumerate}[leftmargin=2em,itemsep=0.25em]
\item \textbf{Initialization.} Set $\ubar_i^{(0)}=1$, $\Abar_j^{(0)}=1$.
\item \textbf{Cost index.} Compute the bilateral cost index
      $\tijk^{\obs}$ once from
      $A^{\obs}_{kl}=(\tkl^{\obs})^{-\theta}$ and
      \eqref{eq:tau_series}.
\item \textbf{Update $\ubar_i$.} For each $i$,
\begin{equation}\label{eq:update_u}
  \ubar_i^{(k+1),\theta}
  \;=\;
  \frac{(l_i^{R,\obs})^{1-\beta\theta}}
       {\chi^{(k)}\displaystyle\sum_{j}(\tijk^{\obs})^{-\theta}
          \Abar_j^{(k),\theta}(l_j^{F,\obs})^{\alpha\theta}}.
\end{equation}
\item \textbf{Update $\Abar_j$.} For each $j$,
\begin{equation}\label{eq:update_A}
  \Abar_j^{(k+1),\theta}
  \;=\;
  \frac{(l_j^{F,\obs})^{1-\alpha\theta}}
       {\chi^{(k)}\displaystyle\sum_{i}(\tijk^{\obs})^{-\theta}
          \ubar_i^{(k+1),\theta}(l_i^{R,\obs})^{\beta\theta}}.
\end{equation}
\item \textbf{Normalization.} Rescale so that
      $\prod_i\ubar_i^{(k+1)}=1$ and $\prod_j\Abar_j^{(k+1)}=1$;
      update $\chi^{(k+1)}$ accordingly from
      \eqref{eq:chi_def}.
\item \textbf{Convergence.} Repeat steps (3)--(5) until
      $\max_i|\ubar_i^{(k+1)}/\ubar_i^{(k)}-1|$ and
      $\max_j|\Abar_j^{(k+1)}/\Abar_j^{(k)}-1|$ fall below a
      tolerance (e.g.\ $10^{-8}$).
\end{enumerate}

\paragraph{Well-posedness.}
Under the uniqueness conditions \eqref{eq:AAunique} of
Proposition~1(b) in \cite{AllenArkolakis2022} and
Assumption~\ref{ass:specrad} applied to $A^{\obs}$, the level
equilibrium system \eqref{eq:eq_R}--\eqref{eq:eq_F} admits a unique
strictly positive solution, and the iteration
\eqref{eq:update_u}--\eqref{eq:update_A} is a \emph{monotone}
operator on the positive orthant (up to the $\chi$ normalization).
We do not establish global contractivity in closed form; instead,
we verify numerically that the spectral radius of the Jacobian of
the fixed-point map at each iterate is strictly less than one, a
sufficient condition for local linear convergence, and check that
the solution is invariant to the initialization.  In our empirical
implementation the algorithm converges in fewer than 50 iterations
from a flat initialization for the square 3\,km grid.

\paragraph{Output of the inversion.}
The inversion returns a baseline profile
$\{(\ubar_i^{\star},\Abar_j^{\star})\}_{i,j\in\N}$ and, as a
by-product, the baseline welfare index $\Wbar^{\star}$.  Together
with the observed link primitives, these fully pin down the baseline
equilibrium and are held fixed in counterfactual analysis.

\subsection{Counterfactuals: Exact-Hat Algebra (Main)}\label{ssec:hat}

Counterfactual analysis is implemented in the main text using the
exact-hat algebra of \cite{AllenArkolakis2022}.  The advantages
relative to a levels computation are well known: the hat system
cancels the unobserved fundamentals $(\ubar_i,\Abar_j)$, is less
sensitive to normalization, and is numerically more stable when
fundamentals are close to the boundary of the allocation.  A levels
counterpart is described in Appendix~\ref{app:levels} and reported
as a robustness check.

\paragraph{Counterfactual policy.}
A tidal-lane policy reallocates directed lane capacity on a target
set $\E^{\star}\subseteq\E$ from the off-peak to the peak direction,
\begin{equation}\label{eq:cf_lanes}
  n_{kl}^{\cf}=n_{kl}+\Delta_{kl},\quad
  n_{lk}^{\cf}=n_{lk}-\Delta_{kl},\quad
  (k,l)\in\E^{\star},
  \qquad
  \Delta_{kl}\in\bigl[0,\,n_{lk}\bigr],
\end{equation}
with all other links unchanged.  The total corridor capacity
$n_{kl}+n_{lk}$ is preserved; only the directional split moves.

\paragraph{Hat notation.}
For any endogenous variable $x$, let $\hat x\equiv x^{\cf}/x$.  The
model primitives that change by assumption are the lane capacities,
$\hat n_{kl}=n_{kl}^{\cf}/n_{kl}$, on the treated links; $\hat n_{kl}=1$
elsewhere.  The free-flow travel times do not change:
$\hat{\tbar}_{kl}=1$ for all $(k,l)$.  All other hats
$(\hat t_{kl},\hat\xi_{kl},\hat\tijk,\hat L_{ij},\hat l_i^R,\hat
l_j^F,\hat P_k,\hat\Pi_l,\hat\Wbar)$ are to be solved for.

\paragraph{Hat system.}
Dividing the baseline and counterfactual versions of each equilibrium
condition yields the following system.

\emph{(i) Congestion technology.}  From
\eqref{eq:cong_primitive},
\begin{equation}\label{eq:hat_cong}
  \hat t_{kl}
  \;=\;
  \hat\xi_{kl}^{\lambda}\,\hat n_{kl}^{-\lambda}.
\end{equation}

\emph{(ii) Bilateral cost index.}  From \eqref{eq:tau_series},
$\tijk^{-\theta}=[(I-A)^{-1}]_{ij}$.  In hat form,
\begin{equation}\label{eq:hat_tau}
  \hat\tijk^{-\theta}
  \;=\;
  \frac{\bigl[(I-A^{\cf})^{-1}\bigr]_{ij}}
       {\bigl[(I-A)^{-1}\bigr]_{ij}},
  \qquad
  A^{\cf}_{kl}=\tkl^{-\theta}\,\hat t_{kl}^{-\theta}.
\end{equation}

\emph{(iii) Link-traffic gravity.}  From \eqref{eq:xi_gravity},
\begin{equation}\label{eq:hat_xi}
  \hat\xi_{kl}
  \;=\;
  \hat t_{kl}^{-\theta}\,\hat P_k^{-\theta}\,\hat \Pi_l^{-\theta}.
\end{equation}

\emph{(iv) Market access.}  From \eqref{eq:Pi_def}--\eqref{eq:Pj_def},
\begin{align}
  \hat\Pi_i^{-\theta}
  &=\frac{\sum_j (\tijk^{\obs})^{-\theta}\hat\tijk^{-\theta}
              L_j^{F,\obs}\hat l_j^F\,P_j^{\theta}\hat P_j^{\theta}}
             {\Pi_i^{-\theta}},
  \label{eq:hat_Pi}\\
  \hat P_j^{-\theta}
  &=\frac{\sum_i (\tijk^{\obs})^{-\theta}\hat\tijk^{-\theta}
              L_i^{R,\obs}\hat l_i^R\,\Pi_i^{\theta}\hat\Pi_i^{\theta}}
             {P_j^{-\theta}}.
  \label{eq:hat_P}
\end{align}

\emph{(v) Share system.}  From \eqref{eq:eq_R}--\eqref{eq:eq_F},
\begin{align}
  \hat l_i^{R,\,1-\beta\theta}
  &=\hat\chi\cdot
    \frac{\sum_j (\tijk^{\obs})^{-\theta}\hat\tijk^{-\theta}
             \Abar_j^{\star,\theta}(l_j^{F,\obs})^{\alpha\theta}
             \hat l_j^{F,\alpha\theta}}
         {\sum_j (\tijk^{\obs})^{-\theta}\Abar_j^{\star,\theta}
             (l_j^{F,\obs})^{\alpha\theta}},
  \label{eq:hat_lR}\\
  \hat l_j^{F,\,1-\alpha\theta}
  &=\hat\chi\cdot
    \frac{\sum_i (\tijk^{\obs})^{-\theta}\hat\tijk^{-\theta}
             \ubar_i^{\star,\theta}(l_i^{R,\obs})^{\beta\theta}
             \hat l_i^{R,\beta\theta}}
         {\sum_i (\tijk^{\obs})^{-\theta}\ubar_i^{\star,\theta}
             (l_i^{R,\obs})^{\beta\theta}},
  \label{eq:hat_lF}
\end{align}
with $\hat\chi=\hat\Wbar^{-\theta}$ fixed by the adding-up
constraints $\sum_i l_i^{R,\obs}\hat l_i^R=1$ and
$\sum_j l_j^{F,\obs}\hat l_j^F=1$.

\begin{remark}[Why fundamentals appear in the hat equations]
Relative to the trade-style exact-hat derivations where
fundamentals drop out entirely, our share system retains
$(\ubar_i^{\star},\Abar_j^{\star})$ because $\alpha\theta$ and
$\beta\theta$ are not $0$ and the baseline fundamentals enter through
the nonlinear externality terms.  Since $(\ubar^{\star},\Abar^{\star})$
are recovered analytically in Section~\ref{ssec:invert}, this causes
no additional computational burden.
\end{remark}

\paragraph{Algorithm (nested fixed point).}

\begin{enumerate}[leftmargin=2em,itemsep=0.25em]
\item \textbf{Inputs.} Baseline
      $\{\tkl^{\obs},L_i^{R,\obs},L_j^{F,\obs},\tijk^{\obs},\xi_{kl}^{\obs},
      P_k,\Pi_l,\ubar_i^{\star},\Abar_j^{\star}\}$ and counterfactual
      lane capacities $\{\hat n_{kl}\}$.
\item \textbf{Outer loop (travel times).} Initialize
      $\hat t_{kl}=1$ for all $(k,l)\in\E$.
\item \textbf{Bilateral cost update.} Given $\{\hat t_{kl}\}$, form
      $A^{\cf}_{kl}=\tkl^{-\theta}\hat t_{kl}^{-\theta}$ and compute
      $\hat\tijk^{-\theta}=[(I-A^{\cf})^{-1}]_{ij}/[(I-A)^{-1}]_{ij}$
      via a single sparse-matrix solve.  This inversion is performed
      once per outer iteration.
\item \textbf{Inner loop (shares).} Given $\{\hat t_{kl},\hat\tijk\}$,
      iterate on \eqref{eq:hat_xi}, \eqref{eq:hat_Pi}--\eqref{eq:hat_P},
      and \eqref{eq:hat_lR}--\eqref{eq:hat_lF} until
      $(\hat\xi_{kl},\hat P_k,\hat\Pi_l,\hat l_i^R,\hat l_j^F)$
      converge (tolerance $10^{-8}$), under the normalizations
      $\sum_i L_i^{R,\obs}\hat l_i^R=\Lbar$,
      $\sum_j L_j^{F,\obs}\hat l_j^F=\Lbar$.
\item \textbf{Outer update.} With inner-loop $\hat\xi_{kl}$, update
      $\hat t_{kl}$ from \eqref{eq:hat_cong}.
\item \textbf{Convergence.} Repeat (3)--(5) until $\hat t$ stops moving
      (tolerance $10^{-6}$).  Terminal
      $(\hat t,\hat\xi,\hat\tau,\hat l^R,\hat l^F,\hat\Wbar)$ is the
      counterfactual equilibrium in proportional changes.
\end{enumerate}

\paragraph{Welfare and distributional metrics.}
Aggregate welfare in proportional changes is
\begin{equation}\label{eq:What}
  \hat\Wbar
  \;=\;\left(\frac{\sum_{ij}(\tijk^{\obs})^{-\theta}\hat\tijk^{-\theta}
                   u_i^{\obs,\theta}\hat u_i^{\theta}
                   w_j^{\obs,\theta}\hat w_j^{\theta}}
                  {\sum_{ij}(\tijk^{\obs})^{-\theta}
                   u_i^{\obs,\theta}w_j^{\obs,\theta}}\right)^{\!1/\theta},
\end{equation}
with $\hat u_i=\hat l_i^{R,\beta}$ and $\hat w_j=\hat l_j^{F,\alpha}$.
Distributional metrics include the within- and between-grid variance
of $\hat l_i^R$ and $\hat l_j^F$, and the fraction of
directional grid pairs whose $\hat\tijk$ moves in the peak-favoring
direction.

\paragraph{Levels counterfactual (Appendix~\ref{app:levels}).}
For robustness we also solve the counterfactual in levels,
re-using the inverted fundamentals
$(\ubar_i^{\star},\Abar_j^{\star})$ and re-solving the equilibrium
system \eqref{eq:eq_R}--\eqref{eq:eq_F} with the counterfactual
$\{\tkl^{\cf},\tijk^{\cf}\}$.  The two solutions should deliver
identical allocations up to numerical tolerance and reporting the
difference is a standard diagnostic.

\subsection{Counterfactual Scenarios}\label{ssec:cf}

The exact-hat solver is applied to four counterfactual scenarios.
All four are solved on the square 3\,km grid at the AM peak with
$(\theta,\alpha,\beta,\lambda)$ held at their baseline calibration;
the PM peak, alternative grid systems, and alternative calibrations
enter as robustness rows.

\paragraph{CF1: Directional vs.\ symmetric commuting costs.}
Replace the observed directional link times with their corridor-level
averages,
\[
  \tkl^{\mathrm{sym}}\;=\;t_{lk}^{\mathrm{sym}}
  \;=\;\tfrac{1}{2}\bigl(\tkl^{\obs}+t_{lk}^{\obs}\bigr),
\]
while holding lane capacities and fundamentals fixed.  This isolates
the model contribution of directional asymmetry holding all other
primitives fixed.  We report the change in $\Wbar$, the cross-grid
distribution of $\hat l_i^R$ and $\hat l_j^F$, and the induced
redistribution of link traffic.

\paragraph{CF2: Tidal-lane reallocation on the identified corridor
set $\E^\star$.}
Apply \eqref{eq:cf_lanes} to the set of directed links
identified in the descriptive analysis as persistent tidal
candidates.  The baseline reallocation is
$\Delta_{kl}=1$ (one lane moved from off-peak to peak direction)
on all $(k,l)\in\E^\star$, with the constraint that the off-peak
direction retains at least one lane.  Sensitivity rows vary
$\Delta_{kl}\in\{0.5,1,2\}$ and the size of $\E^\star$ (top 5\%,
10\%, 20\% most-asymmetric corridors).  Headline outputs:
$\hat\Wbar$, the mean/median $\hat\tijk$ change by residence--workplace
pair, the share of grids that gain $\hat l_i^R>1$, and the
distributional Lorenz curves on $\hat\tijk$.

\paragraph{CF3: Global road expansion.}
Increase $n_{kl}$ uniformly on all links by a factor of $1.1$,
$1.25$, or $1.5$.  This benchmark comparison sets the cost of
undifferentiated capacity against the targeted reallocation in CF2,
and delivers evidence on the ``fundamental law of road congestion''
inside the model.  The $\lambda$-sensitivity is the central
robustness axis.  At low $\lambda$, uniform expansion produces only
modest welfare gains, since congested travel time is nearly
inelastic in directed capacity.  At high $\lambda$, the welfare
response to expansion becomes first-order, and CF2 retains the
interpretation of a second-best reallocation of a fixed lane budget.

\paragraph{CF4: Welfare and distributional comparisons.}
A consolidated comparison across CF1--CF3.  We report
\[
  \mathrm{WelfareGain}\;=\;\ln\hat\Wbar,
  \qquad
  \mathrm{Gini}(\hat l^R),\ \mathrm{Gini}(\hat l^F),
  \qquad
  \mathrm{Pct.\;gainers}\;=\;\frac{1}{N}\sum_i\mathbf 1\{\hat l_i^R>1\},
\]
along with the cross-scenario correlation of changes at the grid
level.  This subsection is intentionally lean on specifics; the
final quantitative results are reported in the main empirical
section of the paper.

\paragraph{Robustness grid.}
For each CF we re-run the solver on (i) the PM peak, (ii) the
hexagonal grid, (iii) $\theta\in\{4,6.83,9\}$, (iv)
$\lambda\in\{0.02,0.05,0.10\}$, and (v) the IV $\hat\lambda$ from
Appendix~\ref{app:iv} where applicable.  These deliver a
multi-dimensional robustness array that accompanies the main
counterfactual table.

%====================================================================
\appendix
\section{Mathematical Appendix}\label{app:math}
%====================================================================

\subsection{Geometric Series Representation of $\tau_{ij}^{-\theta}$}\label{app:tau}

\begin{proof}[Proof of Proposition~\ref{prop:tau_series}]
Fix $(i,j)\in\N^2$.  Every feasible route $r\in\R_{ij}$ has finite
length $K(r)\ge 0$, with $r_0=i$ and $r_{K(r)}=j$, and every edge
$(r_{m-1},r_m)$ belongs to $\E$.  Hence
\[
  \R_{ij}
  \;=\;\bigcup_{K=0}^{\infty}\R_{ij}(K),
  \qquad
  \R_{ij}(K)=\bigl\{\text{directed walks of length $K$ from $i$ to $j$}\bigr\}.
\]
Splitting the sum in \eqref{eq:tau_def} by walk length gives
\[
  \tijk^{-\theta}
  \;=\;\sum_{K=0}^{\infty}\sum_{r\in\R_{ij}(K)}
       \prod_{m=1}^{K}t_{r_{m-1},r_m}^{-\theta}
  \;=\;\sum_{K=0}^{\infty}\sum_{r\in\R_{ij}(K)}
       \prod_{m=1}^{K}a_{r_{m-1},r_m},
\]
where the second equality uses \eqref{eq:Adef}.  By the
combinatorial interpretation of matrix powers,
\[
  \sum_{r\in\R_{ij}(K)}\prod_{m=1}^{K}a_{r_{m-1},r_m}
  \;=\;\bigl(A^{K}\bigr)_{ij}.
\]
Therefore
\[
  \tijk^{-\theta}=\sum_{K=0}^{\infty}\bigl(A^{K}\bigr)_{ij}.
\]
Under Assumption~\ref{ass:specrad}, $\spr(A)<1$, so
$\sum_{K\ge0}A^{K}$ converges absolutely and equals $(I-A)^{-1}$
\cite[see, e.g.,]{AllenArkolakis2022}.  Extracting the $(i,j)$ entry,
$\tijk^{-\theta}=\bigl[(I-A)^{-1}\bigr]_{ij}=b_{ij}$, and
$\tijk=b_{ij}^{-1/\theta}$.
\end{proof}

\begin{remark}[Length-zero walks]
The $K=0$ term contributes $(A^{0})_{ij}=\delta_{ij}$.  This corresponds
to the convention that a worker choosing $i=j$ incurs zero link
cost; the route $(i)$ alone is formally feasible and enters the
gravity sum.  Alternative conventions that exclude within-grid
commuting set this term to zero and require small modifications to
the marginal-population system \eqref{eq:marketclearing}.
\end{remark}

\subsection{Max-of-Fr\'echet and the Welfare Index}\label{app:welfare}

\paragraph{Setup.}
For each triple $(i,j,r)$ the agent has $V_{ij,r}=\omega_{ij,r}
\varepsilon_{ij,r}$ with $\omega_{ij,r}$ as in \eqref{eq:omega_def}
and $\varepsilon_{ij,r}$ i.i.d.\ Fr\'echet$(\theta)$ with unit scale.
Standard properties give:
\[
  \Pr(\varepsilon_{ij,r}\le x)=\exp(-x^{-\theta})\ \ (x>0),
\]
and for any scale $\omega>0$, $\omega\varepsilon$ is itself Fr\'echet
with shape $\theta$ and scale $\omega$:
$\Pr(\omega\varepsilon\le y)=\exp(-(y/\omega)^{-\theta})$.

\paragraph{Maximum of independent Fr\'echet random variables.}
Let $M=\max_{(i,j,r)}V_{ij,r}$.  By independence across $(i,j,r)$,
\begin{align*}
  \Pr(M\le y)
  &=\prod_{(i,j,r)}\Pr(V_{ij,r}\le y)
   =\prod_{(i,j,r)}\exp\bigl(-(y/\omega_{ij,r})^{-\theta}\bigr)\\
  &=\exp\!\left(-y^{-\theta}\sum_{(i,j,r)}\omega_{ij,r}^{\theta}\right)
   =\exp\bigl(-y^{-\theta}S\bigr),
\end{align*}
with $S\equiv\sum_{(i,j,r)}\omega_{ij,r}^{\theta}$, as claimed.

\paragraph{$S$ collapses to a bilateral index.}
Substitute $\omega_{ij,r}=u_i w_j\prod_{m=1}^{K(r)}t_{r_{m-1},r_m}^{-1}$:
\[
  \sum_{r\in\R_{ij}}\omega_{ij,r}^{\theta}
  =u_i^{\theta}w_j^{\theta}\,\sum_{r\in\R_{ij}}
     \prod_{m=1}^{K(r)}t_{r_{m-1},r_m}^{-\theta}
  =u_i^{\theta}w_j^{\theta}\,\tijk^{-\theta}.
\]
Summing over $(i,j)$ delivers \eqref{eq:S_def}:
$S=\sum_{i,j}u_i^{\theta}w_j^{\theta}\tijk^{-\theta}$.

\paragraph{Welfare index.}
$M$ is Fr\'echet with shape $\theta$ and scale $S^{1/\theta}$, so
$\mathbb E[M]=\Gamma(1-1/\theta)S^{1/\theta}$.  Absorbing the
gamma constant into the welfare index, we define
$\Wbar\equiv S^{1/\theta}$, consistent with \eqref{eq:W_def}.

\paragraph{Gravity from the distribution of $M$.}
To recover \eqref{eq:gravity}, note that for Fr\'echet-max systems,
the probability that option $(i,j,r)$ attains the maximum is its
scale raised to $\theta$ divided by the total scale, i.e.\
$\omega_{ij,r}^{\theta}/S$.  Summing over routes,
$\Pr(i,j)=u_i^{\theta}w_j^{\theta}\tijk^{-\theta}/S$, and
multiplying by $\Lbar$ gives
$L_{ij}=\tijk^{-\theta}u_i^{\theta}w_j^{\theta}\Lbar/\Wbar^{\theta}$.

\subsection{Welfare Derivations}\label{app:welfare2}

This appendix expands the derivations asserted in
\S\ref{ssec:welfare}.

\paragraph{Proof of \eqref{eq:welfare_main}.}
The maximum $M=\max_{(i,j,r)}V_{ij,r}$ of independent Fr\'echet random
variables $V_{ij,r}=\omega_{ij,r}\varepsilon_{ij,r}$ with
$\varepsilon_{ij,r}\sim\text{Fr\'echet}(\theta)$ and scales
$\omega_{ij,r}=u_iw_j/\prod_m t_{r_{m-1},r_m}$ has CDF
\[
  \Pr(M\le y)=\prod_{(i,j,r)}\Pr(\omega_{ij,r}\varepsilon_{ij,r}\le y)
  =\exp\!\Bigl(-y^{-\theta}\sum_{(i,j,r)}\omega_{ij,r}^{\theta}\Bigr)
  =\exp(-y^{-\theta}S),
\]
with $S=\sum_{(i,j)}\tijk^{-\theta}u_i^{\theta}w_j^{\theta}$ by the
definition of $\tijk^{-\theta}$ in \eqref{eq:tau_def}.  This is the
Fr\'echet CDF with shape $\theta$ and scale $S^{1/\theta}=\Wbar$.
Since $\Gamma(1-1/\theta)$ is the expectation of the unit-scale
Fr\'echet, $\mathbb E[M]=\Gamma(1-1/\theta)\Wbar$, which is
\eqref{eq:welfare_main}.

\paragraph{Proof of \eqref{eq:welfare_MA}.}
Using the market-access definition \eqref{eq:Pi_def},
\[
  \Pi_i^{-\theta}
  =\sum_{j\in\N}\tijk^{-\theta}L_j^FP_j^{\theta}
  =\sum_{j\in\N}\tijk^{-\theta}\frac{L_j^F}{L_j^F\Lbar^{-1}\cdot P_j^{-\theta}}
  \cdot (\text{after substituting \eqref{eq:Pj_def}}),
\]
and similar algebra applied to $P_j^{-\theta}$.  After simplification,
$u_i^{\theta}\Pi_i^{-\theta}=\sum_j\tijk^{-\theta}u_i^{\theta}w_j^{\theta}$
up to multiplicative constants that cancel by
\eqref{eq:marketclearing}, delivering
$\Wbar^{\theta}=\sum_iu_i^{\theta}\Pi_i^{-\theta}$.  The workplace
version follows by symmetry.

\paragraph{Proof of \eqref{eq:welfare_constant_across_pairs}.}
Conditional on $(i,j)$ being chosen, the distribution of $V$ is the
conditional Fr\'echet given it attains the maximum, which has the
\emph{same} scale $\Wbar$ as the unconditional maximum (a standard
property of extreme-value distributions).  Hence
$\mathbb E[V\mid(i,j)]=\Gamma(1-1/\theta)\Wbar=\mathbb W$ for all
chosen $(i,j)$.

\paragraph{Derivation of \eqref{eq:welfare_decomp}.}
Log-differentiating $\Wbar^{\theta}=\sum_{(i,j)}\tijk^{-\theta}
u_i^{\theta}w_j^{\theta}$ and normalizing by $\Wbar^{\theta}\Lbar$,
\[
  \theta\,\mathrm d\ln\Wbar
  \;=\;\sum_{(i,j)}\frac{\tijk^{-\theta}u_i^{\theta}w_j^{\theta}}{\Wbar^{\theta}}
       \bigl(\theta\,\mathrm d\ln u_i+\theta\,\mathrm d\ln w_j
        -\theta\,\mathrm d\ln\tijk\bigr).
\]
By the gravity equation \eqref{eq:gravity},
$\tijk^{-\theta}u_i^{\theta}w_j^{\theta}/\Wbar^{\theta}
=L_{ij}/\Lbar=\Lambda_{ij}$, and summing out along $j$ (resp.\ $i$)
yields $\sum_j\Lambda_{ij}=\Lambda_i^R$ (resp.\
$\sum_i\Lambda_{ij}=\Lambda_j^F$).  Cancelling $\theta$ and using
$\mathrm d\ln\mathbb W=\mathrm d\ln\Wbar$ gives
\eqref{eq:welfare_decomp}.

\paragraph{Derivation of the hat-welfare formula
\eqref{eq:welfare_hat}.}
Write the counterfactual welfare scale as
\[
  \Wbar'^{\theta}
  =\sum_{(i,j)}\tau_{ij}'^{-\theta}u_i'^{\theta}w_j'^{\theta}
  =\sum_{(i,j)}\hat{\tau}_{ij}^{-\theta}\hat u_i^{\theta}\hat w_j^{\theta}
   \cdot\tijk^{-\theta}u_i^{\theta}w_j^{\theta}.
\]
Divide both sides by $\Wbar^{\theta}$ and substitute
$\tijk^{-\theta}u_i^{\theta}w_j^{\theta}/\Wbar^{\theta}=\pi_{ij}^{\obs}
=L_{ij}^{\obs}/\Lbar$ from the baseline gravity equation to obtain
\[
  \hat{\Wbar}^{\theta}
  =\sum_{(i,j)}\pi_{ij}^{\obs}\,\hat{\tau}_{ij}^{-\theta}
   \,\hat u_i^{\theta}\,\hat w_j^{\theta},
\]
which is \eqref{eq:welfare_hat}.  Equation \eqref{eq:hat_uw} follows
from the spillover equations \eqref{eq:fundamentals} by direct
computation: $\hat u_i=(L_i^{R\prime}/L_i^R)^{\beta}=\hat
l_i^{R,\beta}$ since the ratios
$\Lbar'/\Lbar=1$ under the share normalisation used in the
counterfactual block, and likewise for $\hat w_j$.

\paragraph{Proof of Proposition~\ref{prop:welfare_bounds}.}
Start from \eqref{eq:welfare_hat} and partition the outer sum by
residence location:
\[
  \hat{\Wbar}^{\theta}
  =\sum_{i}\sum_{j}\pi_{ij}^{\obs}\hat{\tau}_{ij}^{-\theta}
   \hat u_i^{\theta}\hat w_j^{\theta}
  =\sum_{i}\hat u_i^{\theta}\underbrace{\sum_{j}\pi_{ij}^{\obs}
   \hat{\tau}_{ij}^{-\theta}\hat w_j^{\theta}}_{=\Lambda_i^R\cdot
   \hat\Pi_i^{-\theta}}
  =\sum_{i}\Lambda_i^R\hat u_i^{\theta}\hat\Pi_i^{-\theta},
\]
where the inner bracket uses
$\pi_{ij}^{\obs}=\Lambda_{ij}=\Lambda_i^R\cdot\pi_{ij|i}^{\obs}$ with
$\pi_{ij|i}^{\obs}=L_{ij}/L_i^R$ and the definition
$\hat\Pi_i^{-\theta}=\sum_j\pi_{ij|i}^{\obs}\hat{\tau}_{ij}^{-\theta}
\hat w_j^{\theta}$.  The workplace version
$\sum_j\Lambda_j^F\hat w_j^{\theta}\hat P_j^{-\theta}$ is obtained by
partitioning on $j$ first, using
$\pi_{ij}^{\obs}=\Lambda_j^F\pi_{ij|j}^{\obs}$.

\subsection{Link-Traffic Gravity: Proof}\label{app:linkgravity}

\paragraph{Step 1: Link-traversal weights.}
Consider the set of walks from $i$ to $j$ that pass through the
directed edge $(k,l)$ at least once.  Any such walk can be written
(allowing for repetitions) as a concatenation of an $i$-to-$k$ walk,
the single edge $(k,l)$, and an $l$-to-$j$ walk.  The total weight of
such walks, with the same $(a_{kl}$-product) weighting as in
\eqref{eq:tau_def}, equals
\[
  \Bigl(\sum_{r\in\R_{ik}}\prod_{m}t_{r_{m-1},r_m}^{-\theta}\Bigr)\cdot
  \tkl^{-\theta}\cdot
  \Bigl(\sum_{r'\in\R_{lj}}\prod_{m}t_{r'_{m-1},r'_m}^{-\theta}\Bigr)
  \;=\;\tau_{ik}^{-\theta}\,\tkl^{-\theta}\,\tau_{lj}^{-\theta},
\]
by the definition of $\tau$ for each sub-walk.

\paragraph{Step 2: Expected number of traversals.}
For a representative commuter from $(i,j)$, the expected number of
times she traverses $(k,l)$ equals the total weight of walks through
$(k,l)$ divided by the total weight of all walks between $i$ and $j$:
\begin{equation}
  \pi^{kl}_{ij}
  \;=\;\frac{\tau_{ik}^{-\theta}\,\tkl^{-\theta}\,\tau_{lj}^{-\theta}}
            {\tijk^{-\theta}}
  \;=\;\left(\frac{\tijk}{\tau_{ik}\,\tkl\,\tau_{lj}}\right)^{\!\theta},
  \label{eq:piijkl_app}
\end{equation}
which is \eqref{eq:piijkl}.

\paragraph{Step 3: Aggregate link traffic.}
Using gravity in its market-access form \eqref{eq:gravity_MA} and
\eqref{eq:piijkl_app}:
\begin{align*}
  \xi_{kl}
  &=\sum_{i,j}L_{ij}\,\pi^{kl}_{ij}
   =\sum_{i,j}\left[\tijk^{-\theta}\frac{L_i^R}{\Pi_i^{-\theta}}
                       \frac{L_j^F}{P_j^{-\theta}}\right]
             \left[\frac{\tau_{ik}^{-\theta}\,\tkl^{-\theta}\,\tau_{lj}^{-\theta}}
                        {\tijk^{-\theta}}\right]\\
  &=\tkl^{-\theta}
    \sum_{i,j}\tau_{ik}^{-\theta}\,\tau_{lj}^{-\theta}\,
              \frac{L_i^R}{\Pi_i^{-\theta}}\,
              \frac{L_j^F}{P_j^{-\theta}}\\
  &=\tkl^{-\theta}\cdot
    \underbrace{\left(\sum_i\tau_{ik}^{-\theta}\,L_i^R\,\Pi_i^{\theta}\right)}_{=P_k^{-\theta}
     \text{ by \eqref{eq:Pj_def}}}\cdot
    \underbrace{\left(\sum_j\tau_{lj}^{-\theta}\,L_j^F\,P_j^{\theta}\right)}_{=\Pi_l^{-\theta}
     \text{ by \eqref{eq:Pi_def}}}.
\end{align*}
Hence $\xi_{kl}=\tkl^{-\theta}P_k^{-\theta}\Pi_l^{-\theta}$, which is
\eqref{eq:xi_gravity}.  This completes the proof of
Proposition~\ref{prop:linkgravity}.\hfill$\square$

\subsection{Equilibrium System in Population Shares}\label{app:eqshares}

We derive \eqref{eq:eq_R}--\eqref{eq:eq_F} step by step.

\paragraph{Step 1: Substitute fundamentals.}
From the gravity equation \eqref{eq:gravity} and the fundamentals
\eqref{eq:fundamentals}--\eqref{eq:wage},
\begin{align*}
  L_{ij}
  &=\tijk^{-\theta}\,u_i^{\theta}\,w_j^{\theta}\,\frac{\Lbar}{\Wbar^{\theta}}\\
  &=\tijk^{-\theta}\,\ubar_i^{\theta}(L_i^R)^{\beta\theta}\,
       \Abar_j^{\theta}(L_j^F)^{\alpha\theta}\,\frac{\Lbar}{\Wbar^{\theta}}.
\end{align*}

\paragraph{Step 2: Sum over $j$.}
\[
  L_i^R=\sum_j L_{ij}
  =\ubar_i^{\theta}(L_i^R)^{\beta\theta}\,\frac{\Lbar}{\Wbar^{\theta}}\,
     \sum_j\tijk^{-\theta}\Abar_j^{\theta}(L_j^F)^{\alpha\theta}.
\]
Re-arranging,
\[
  (L_i^R)^{1-\beta\theta}
  =\frac{\Lbar}{\Wbar^{\theta}}\,\ubar_i^{\theta}\,
     \sum_j\tijk^{-\theta}\Abar_j^{\theta}(L_j^F)^{\alpha\theta}.
  \tag{$*$}
\]

\paragraph{Step 3: Convert to shares.}
Using $l_i^R=L_i^R/\Lbar$ and $l_j^F=L_j^F/\Lbar$,
$L_i^R=\Lbar l_i^R$ and $L_j^F=\Lbar l_j^F$, so
$(L_i^R)^{1-\beta\theta}=\Lbar^{1-\beta\theta}(l_i^R)^{1-\beta\theta}$
and $(L_j^F)^{\alpha\theta}=\Lbar^{\alpha\theta}(l_j^F)^{\alpha\theta}$.
Substituting into ($*$),
\[
  \Lbar^{1-\beta\theta}(l_i^R)^{1-\beta\theta}
  =\frac{\Lbar}{\Wbar^{\theta}}\,\ubar_i^{\theta}\,
     \sum_j\tijk^{-\theta}\Abar_j^{\theta}\Lbar^{\alpha\theta}
     (l_j^F)^{\alpha\theta},
\]
and dividing both sides by $\Lbar^{1-\beta\theta}$ gives
\[
  (l_i^R)^{1-\beta\theta}
  =\frac{\Lbar^{\alpha\theta+\beta\theta}}{\Wbar^{\theta}}\,
     \ubar_i^{\theta}\,
     \sum_j\tijk^{-\theta}\Abar_j^{\theta}(l_j^F)^{\alpha\theta}.
\]
Defining $\chi\equiv\Lbar^{\theta(\alpha+\beta)}/\Wbar^{\theta}$ yields
\eqref{eq:eq_R}.

\paragraph{Step 4: Symmetric argument for $l_j^F$.}
Summing $L_{ij}$ over $i$ with the gravity equation and repeating
the substitutions delivers \eqref{eq:eq_F}.  The summation index in
$\sum_i\tijk^{-\theta}\ubar_i^{\theta}(l_i^R)^{\beta\theta}$ uses
$\tijk^{-\theta}$ (going from $i$ to $j$, fixed $j$); this is the
correct directional cost for the firm-side equation.

\subsection{Exact-Hat Algebra: Full System}\label{app:hat}

This appendix collects all hat equations in one place for reference.
Let $\hat x=x^{\cf}/x$ for any endogenous variable $x$.  The
exogenous policy shock is $\hat n_{kl}$.  Free-flow times are
unchanged, $\hat{\tbar}_{kl}=1$.  The baseline primitives
$(L^{R,\obs},L^{F,\obs},\tkl^{\obs},\tijk^{\obs},\xi_{kl}^{\obs},
P_k,\Pi_l,\ubar^{\star},\Abar^{\star},\Wbar^{\star})$ are known from
the data and the inversion.

\paragraph{Endogenous hats to solve for.}
\[
  \hat t_{kl},\quad\hat\xi_{kl},\quad\hat\tijk,\quad\hat P_k,
  \quad\hat\Pi_l,\quad\hat l_i^R,\quad\hat l_j^F,\quad\hat\Wbar.
\]

\paragraph{Equations.}
\begin{align}
  \hat t_{kl}
  &=\hat\xi_{kl}^{\lambda}\hat n_{kl}^{-\lambda},
  \label{eq:hat_cong_app}\\[2pt]
  \hat\tijk^{-\theta}
  &=\frac{[(I-A^{\cf})^{-1}]_{ij}}{[(I-A^{\obs})^{-1}]_{ij}},
    \quad A^{\cf}_{kl}=\tkl^{\obs,-\theta}\hat t_{kl}^{-\theta},
  \label{eq:hat_tau_app}\\[2pt]
  \hat\xi_{kl}
  &=\hat t_{kl}^{-\theta}\hat P_k^{-\theta}\hat\Pi_l^{-\theta},
  \label{eq:hat_xi_app}\\[2pt]
  \hat\Pi_i^{-\theta}
  &=\frac{\sum_j (\tijk^{\obs})^{-\theta}\hat\tijk^{-\theta}
                 L_j^{F,\obs}\hat l_j^F P_j^{\theta}\hat P_j^{\theta}}
             {\Pi_i^{-\theta}},
  \label{eq:hat_Pi_app}\\[2pt]
  \hat P_j^{-\theta}
  &=\frac{\sum_i (\tijk^{\obs})^{-\theta}\hat\tijk^{-\theta}
                 L_i^{R,\obs}\hat l_i^R \Pi_i^{\theta}\hat\Pi_i^{\theta}}
             {P_j^{-\theta}},
  \label{eq:hat_Pj_app}\\[2pt]
  \hat l_i^{R,1-\beta\theta}
  &=\hat\chi\cdot
    \frac{\sum_j (\tijk^{\obs})^{-\theta}\hat\tijk^{-\theta}
                 \Abar_j^{\star,\theta}(l_j^{F,\obs})^{\alpha\theta}
                 \hat l_j^{F,\alpha\theta}}
         {\sum_j (\tijk^{\obs})^{-\theta}
                 \Abar_j^{\star,\theta}(l_j^{F,\obs})^{\alpha\theta}},
  \label{eq:hat_lR_app}\\[2pt]
  \hat l_j^{F,1-\alpha\theta}
  &=\hat\chi\cdot
    \frac{\sum_i (\tijk^{\obs})^{-\theta}\hat\tijk^{-\theta}
                 \ubar_i^{\star,\theta}(l_i^{R,\obs})^{\beta\theta}
                 \hat l_i^{R,\beta\theta}}
         {\sum_i (\tijk^{\obs})^{-\theta}
                 \ubar_i^{\star,\theta}(l_i^{R,\obs})^{\beta\theta}},
  \label{eq:hat_lF_app}\\[2pt]
  \hat\Wbar^{\theta}
  &=\frac{\sum_{ij}(\tijk^{\obs})^{-\theta}\hat\tijk^{-\theta}
                   u_i^{\obs,\theta}\hat u_i^{\theta}
                   w_j^{\obs,\theta}\hat w_j^{\theta}}
          {\sum_{ij}(\tijk^{\obs})^{-\theta}
                   u_i^{\obs,\theta}w_j^{\obs,\theta}},
  \label{eq:hat_W_app}\\[2pt]
  \hat u_i&=\hat l_i^{R,\beta},\qquad
  \hat w_j=\hat l_j^{F,\alpha},
  \label{eq:hat_uw_app}\\[2pt]
  \hat\chi&=\hat\Wbar^{-\theta},
  \qquad
  \sum_i l_i^{R,\obs}\hat l_i^R=1,
  \qquad
  \sum_j l_j^{F,\obs}\hat l_j^F=1.
  \label{eq:hat_norm_app}
\end{align}

\paragraph{Solution procedure.}
The nested fixed-point solver described in Section~\ref{ssec:hat}
computes the inner fixed point
\eqref{eq:hat_cong_app}--\eqref{eq:hat_Pj_app} for given
$(\hat l^R,\hat l^F)$, then updates the outer variables using
\eqref{eq:hat_lR_app}--\eqref{eq:hat_lF_app} and the normalization
\eqref{eq:hat_norm_app}.  Convergence diagnostics and alternative
updating schemes (Anderson acceleration, damped Newton) are
implemented in the companion code.

\subsection{Counterfactuals in Levels}\label{app:levels}

As a robustness check for the exact-hat solver, we also compute
counterfactuals in levels, holding the inverted fundamentals
$(\ubar_i^{\star},\Abar_j^{\star})$ and the parameters
$(\theta,\alpha,\beta,\lambda)$ fixed and re-solving
Definition~\ref{def:eq}.

\paragraph{Inputs.}
$(\ubar_i^{\star},\Abar_j^{\star},\theta,\alpha,\beta,\lambda)$;
counterfactual link primitives
$\{(\ell_{kl},v_{kl}^{0},n_{kl}^{\cf})\}_{(k,l)\in\E}$.

\paragraph{Algorithm.}
\begin{enumerate}[leftmargin=2em,itemsep=0.2em]
\item \textbf{Initialize} $t_{kl}^{(0)}=\tbar_{kl}$ for all $(k,l)$.
\item \textbf{Step A (cost index).} Compute
      $A^{(k)}_{kl}=t_{kl}^{(k),-\theta}$, $\tijk^{(k)}=
      [(I-A^{(k)})^{-1}]_{ij}^{-1/\theta}$.
\item \textbf{Step B (shares).} Solve
      \eqref{eq:eq_R}--\eqref{eq:eq_F} for $(l_i^R,l_j^F)$ given
      $\tijk^{(k)}$ and $\ubar^{\star},\Abar^{\star}$ by the same
      iterative fixed-point algorithm as the inversion.
\item \textbf{Step C (market access).} Given $(l^R,l^F)$ and
      $\tijk^{(k)}$, compute $(P_k,\Pi_l)$ from the linear system
      \eqref{eq:Pi_def}--\eqref{eq:Pj_def} (fixed-point iteration).
\item \textbf{Step D (traffic and travel time).} Compute
      $\xi_{kl}^{(k)}=t_{kl}^{(k),-\theta}P_k^{-\theta}\Pi_l^{-\theta}$
      and update
      $t_{kl}^{(k+1)}=\tbar_{kl}(\xi_{kl}^{(k)}/n_{kl}^{\cf})^{\lambda}$.
\item \textbf{Convergence.} Repeat A--D until
      $\max_{kl}|t_{kl}^{(k+1)}/t_{kl}^{(k)}-1|<10^{-6}$.
\end{enumerate}

\paragraph{Consistency check.}
Multiplying the level counterfactual outputs by the inverse of the
baseline levels should reproduce the exact-hat solution to numerical
tolerance.  The maximum pointwise deviation across
$(\hat t_{kl},\hat\xi_{kl},\hat l_i^R,\hat l_j^F)$ is reported in the
robustness table.

\subsection{Numerical Checks of Equilibrium Uniqueness}\label{app:uniqueness_numeric}

On grid points where the Allen-Arkolakis sufficient condition
\eqref{eq:AAunique} fails, we do not claim uniqueness in closed form.
This appendix records the three numerical diagnostics referenced in
Section~\ref{ssec:eq} and reported for every specification on the
$(\alpha,\beta,\theta,\lambda)$ grid.

\paragraph{Check 1: Multi-start invariance.}
For each parameter configuration, the inversion solver
\eqref{eq:update_u}--\eqref{eq:update_A} and the congestion
fixed-point \eqref{eq:hat_cong_app}--\eqref{eq:hat_Pj_app} are each
run from 50 randomised initialisations drawn from
$\log\ubar_i^{(0)},\log\Abar_j^{(0)}\sim\mathcal U[-1,1]$ and
$\log t_{kl}^{(0)}\sim\mathcal U[\log\tbar_{kl}-0.5,\log\tbar_{kl}+0.5]$.
The reported solution is the solution at which
$\max_{i,j,kl}|x^{(\text{start }s)}-x^{(\text{start }1)}|<10^{-6}$ for
all $s$; a start that fails this tolerance is flagged.

\paragraph{Check 2: Spectral radius of the equilibrium map.}
At each converged solution, the Jacobian of the equilibrium map (in
the shares representation of Appendix~\ref{app:eqshares}) is
computed by finite differences, and its spectral radius $\rho(J)$ is
reported.  A value $\rho(J)<1$ is a sufficient local condition for
isolated uniqueness at that fixed point.

\paragraph{Check 3: Continuity in parameters.}
For every pair of adjacent grid points
$(\alpha,\beta,\theta,\lambda)$ and
$(\alpha,\beta,\theta,\lambda)'$ with a single coordinate perturbed,
the solution distance
$\|\log x-\log x'\|_\infty$ is reported.  A discontinuous jump in the
solution across an infinitesimal parameter perturbation would be
consistent with a change of equilibrium branch; monotone and
continuous variation across the grid is consistent with a single
branch.

\paragraph{Reporting.}
The three diagnostics are tabulated in the counterfactual robustness
output for every grid point.  In every specification reported in
Section~\ref{ssec:cf}, Check~1 is satisfied at the $10^{-6}$
tolerance, Check~2 returns $\rho(J)$ strictly below one, and
Check~3 returns a solution that varies smoothly across the grid.
No formal claim of uniqueness beyond the reach of
\eqref{eq:AAunique} is made.

\subsection{Gravity Estimation of $\theta$ (Internal)}\label{app:theta}

As a robustness complement to the external calibration $\theta=6.83$,
we report an internal gravity-based estimator.

\paragraph{Estimating equation.}
Taking logs of \eqref{eq:gravity} and absorbing origin- and
destination-specific terms into fixed effects yields
\begin{equation}\label{eq:gravity_log}
  \ln L_{ij}^{\obs}
  \;=\;
  -\theta\,\ln\tijk
  +\mu_i+\nu_j+\varepsilon_{ij},
\end{equation}
where $\mu_i$ absorbs $\theta\ln u_i+$ constants and $\nu_j$
absorbs $\theta\ln w_j+$ constants.  Because $\tijk$ depends on
$\theta$ through the series \eqref{eq:tau_series}, \eqref{eq:gravity_log}
cannot be estimated by OLS in $\tijk$ alone.

\paragraph{Grid-search estimator.}
We implement an outer grid search over
$\theta\in\Theta=\{3,4,\ldots,12\}$ with fine refinement near
the minimum:
\begin{enumerate}[leftmargin=2em,itemsep=0.15em]
  \item For each $\theta\in\Theta$, compute
        $A^{\obs}_{kl}(\theta)=(\tkl^{\obs})^{-\theta}$ and
        $\tijk^{\obs}(\theta)=[(I-A^{\obs}(\theta))^{-1}]_{ij}^{-1/\theta}$.
  \item Estimate the two-way fixed-effects regression
        \eqref{eq:gravity_log} by OLS (or PPML) on the reachability
        sample.  Record the log-likelihood (or $R^2$) and the
        estimated slope $\hat\theta_{\mathrm{slope}}(\theta)$.
  \item Select $\hat\theta$ as the value that best fits
        $\hat\theta_{\mathrm{slope}}(\theta)=\theta$, i.e.\ the
        internal-consistency fixed point.
\end{enumerate}
This is the standard consistency check and is numerically identical
to the profile-likelihood estimator when the moment is quadratic.

\paragraph{Sample.}
The regression is run on the largest weakly connected component of
the directed grid graph, dropping self-loop pairs $i=j$ and any
$(i,j)$ with $L_{ij}^{\obs}=0$.  As an alternative, a PPML
specification on all pairs (including zero flows) is reported as a
robustness row.

\paragraph{Relation to main calibration.}
If $\hat\theta$ falls within $[6,8]$, the external calibration
$\theta=6.83$ is retained for the counterfactual; otherwise we
re-run counterfactuals with $\hat\theta$ and report the comparison.

\subsection{IV Strategy for $\lambda$}\label{app:iv}

The naive OLS estimate from \eqref{eq:lambda_regression} is
mechanically circular: the regressor $\hat\xi_{kl}$ is built from
$\{\tkl^{\obs}\}$, which is itself a transformation of the
dependent variable $v_{kl}^{\obs}$ through
$\tkl^{\obs}=\ell_{kl}/v_{kl}^{\obs}$.  Under the free-flow and
leave-one-out constructions of Section~\ref{ssec:lambdahat} the
circularity is broken, but a residual simultaneity remains: in
equilibrium, route choice allocates commuters to links according to
endogenous travel speeds.  The IV strategy addresses both, by
replacing the endogenous regressor with predicted traffic per lane
from instruments that shift the composition of $n_{kl}$ without
affecting contemporaneous speed shocks on $(k,l)$.

\paragraph{Candidate instruments.}

\emph{(i) Historical road capacity.}  Lane counts that predate the
current distribution of commuting activity shift contemporaneous
traffic per lane through the composition of $n_{kl}$ but should be
orthogonal to contemporaneous speed shocks.  Sources: roadtype
attribute at the first observation year in the raw network, or
archival GIS road layers.

\emph{(ii) Distance to CBD and employment clusters.}  The distance
from link $(k,l)$ to the geographic center of Beijing or to the
nearest top-quintile employment grid is strongly correlated with
commuting demand but is a link-geography primitive not affected by
contemporaneous speed.  A simple instrument is
$z^{\mathrm{dist}}_{kl}=\log(1+d_{kl}^{\mathrm{CBD}})$.

\emph{(iii) Reverse-direction nighttime capacity.}  The nighttime
free-flow speed on the \emph{reverse} direction of the same corridor
serves as a shifter for baseline road quality symmetric to $(k,l)$
but is independent of the AM-peak flow on $(k,l)$ conditional on
lane-bin fixed effects.

\paragraph{First stage.}
Letting $z_{kl}$ collect the instruments, the first stage is
\begin{equation}
  \ln\!\left(\frac{\hat\xi_{kl}}{n_{kl}}\right)
  \;=\;
  \gamma_{0}
  +\gamma_{1}'z_{kl}
  +\eta_{\mathrm{bin}(kl)}
  +u_{kl}.
  \label{eq:iv_first}
\end{equation}
Standard relevance and exclusion requirements apply: the $F$-statistic
on $\gamma_{1}$ should exceed conventional thresholds, and any
direct effect of $z_{kl}$ on contemporaneous speed (not through
traffic per lane) should be absent.

\paragraph{Second stage and identification.}
Using the predicted $\widehat{\ln(\hat\xi_{kl}/n_{kl})}$ from
\eqref{eq:iv_first} in place of the endogenous regressor in
\eqref{eq:lambda_regression} delivers the 2SLS estimator
$\hat\lambda^{\mathrm{IV}}$.  Identification rests on the exclusion
restriction that $z_{kl}$ affects $\ln v_{kl}$ only through
$\ln(\xi_{kl}/n_{kl})$, holding lane-bin fixed effects fixed.  Over-
identification tests are reported when more than one instrument is
used.

\paragraph{Baseline and robustness.}
$\hat\lambda^{\mathrm{IV}}$ is the baseline estimate used in the
main counterfactuals.  The naive-OLS
$\hat\lambda^{\mathrm{OLS,naive}}$ and the free-flow and
leave-one-out OLS variants from Section~\ref{ssec:lambdahat} are
reported in the robustness array; the distance between the IV row
and the naive-OLS row is reported explicitly as a diagnostic of the
mechanical circularity documented there.

\paragraph{Discussion and caveats.}
The instruments in (i)--(iii) are scoped but not yet implemented in
the current version of the empirical appendix.  First-stage
$F$-statistics, over-identification tests, and the point estimate
$\hat\lambda^{\mathrm{IV}}$ are left as
[\textsc{todo: report once IV pipeline finalised}].  Until then,
counterfactuals reported under the IV baseline should be read as
indicative, and the free-flow OLS construction of
Section~\ref{ssec:lambdahat} provides an interim lower-bias
estimate of $\lambda$.

\bibliographystyle{plain}
\begin{thebibliography}{10}
\bibitem{AllenArkolakis2014}
Allen, T., Arkolakis, C. (2014). Trade and the topography of the
spatial economy. \emph{Quarterly Journal of Economics}, 129(3),
1085--1140.

\bibitem{AllenArkolakis2022}
Allen, T., Arkolakis, C. (2022). The welfare effects of transportation
infrastructure improvements. \emph{Review of Economic Studies},
89(6), 2911--2957.

\bibitem{Ahlfeldt2015}
Ahlfeldt, G., Redding, S. J., Sturm, D. M., Wolf, N. (2015). The
economics of density: Evidence from the Berlin Wall.
\emph{Econometrica}, 83(6), 2127--2189.

\bibitem{CombesDurantonGobillon2008}
Combes, P.-P., Duranton, G., Gobillon, L. (2008). Spatial wage
disparities: Sorting matters! \emph{Journal of Urban Economics},
63(2), 723--742.

\bibitem{CombesGobillon2015}
Combes, P.-P., Gobillon, L. (2015). The empirics of agglomeration
economies. In \emph{Handbook of Regional and Urban Economics}, Vol.
5A, ed.\ G.\ Duranton, J.\ V.\ Henderson, and W.\ Strange, 247--348.
Elsevier.

\bibitem{MonteReddingRossiHansberg2018}
Monte, F., Redding, S. J., Rossi-Hansberg, E. (2018). Commuting,
migration, and local employment elasticities. \emph{American Economic
Review}, 108(12), 3855--3890.

\bibitem{HeblichReddingSturm2020}
Heblich, S., Redding, S. J., Sturm, D. M. (2020). The making of the
modern metropolis: Evidence from London. \emph{Quarterly Journal of
Economics}, 135(4), 2059--2133.

\bibitem{DurantonTurner2011}
Duranton, G., Turner, M. A. (2011). The fundamental law of road
congestion: Evidence from US cities. \emph{American Economic Review},
101(6), 2616--2652.
\end{thebibliography}

\end{document}
